\chapter{Bioinformatics: sequence, structure, omics}
\label{vol06:ch11}

\epigraph{It has not escaped our notice that the specific pairing we
have postulated immediately suggests a possible copying mechanism for
the genetic material.}{\citeauthor{paper:v6c11-watson-crick-1953},
\citetitle{paper:v6c11-watson-crick-1953}, p.~737}

\chapmeta{Half-life: medium-life. Archetypes: information, scaling, uncertainty. Exercise target: 30. Project track: simulation. Hazard class: none. Reader path default: standard, with sections explicitly tagged. Named cases: Duke Potti genomic predictor clinical trials.}

\noindent
Biology is the largest measurement programme on Earth. A single bacterial
genome holds a few megabytes; a human genome holds about three gigabytes;
a metagenomic sample of soil holds terabytes; the global archive of
sequenced reads passed an exabyte during the early 2020s
\cite{paper:v6c11-stephens-genomic-2015}. The discipline of bioinformatics
is the engineering of how that measurement programme is stored,
searched, aligned, compared, predicted, integrated, and used to support
decisions about cells, organisms, and patients. Each step is a transformation
on data with a known biological structure and an enormous noise budget,
and the engineer's job is to keep the structure that the biology supplies
visible against the noise that the measurement introduces. This chapter
treats sequence data, sequence-to-structure prediction, genomics and
variant calling, transcriptomics, the broader omics stack, phylogenetics,
machine learning for biology, and the canonical failure mode by which
biological machine learning leaks itself into reporting overdrawn results.

\section{Sequence data: alphabets, alignment, search}\pathtag{core}

\subsection{Alphabets, formats, and quality}

A biological sequence is a string over a small alphabet. DNA and RNA use
four symbols (A, C, G, T or U); protein sequences use the twenty
standard amino acids, sometimes extended by codes for selenocysteine
and pyrrolysine. The alphabet matters in two ways: it sets the maximum
information per symbol ($\log_2 4 = 2$~bits for DNA, $\log_2 20
\approx 4.32$~bits for protein), and it constrains the substitution
statistics that an alignment algorithm uses.

Read data are stored most often as FASTA (header plus sequence) or
FASTQ (header, sequence, separator, per-base quality string)
\cite{paper:v6c11-cock-fasta-2010}. The Phred quality score $Q$ encodes
an estimated base-call error probability $p$ as
\[
Q = -10 \log_{10} p,
\]
so $Q = 30$ means a one-in-a-thousand error rate at that position.
Quality scores are encoded as ASCII characters offset by 33 (Sanger
encoding) or 64 (early Illumina); modern pipelines assume Sanger
unless told otherwise. A quality filter trims reads at the position
where the Phred score falls below a threshold (commonly $Q=20$),
because the downstream cost of an erroneous base in a low-quality
position is greater than the cost of throwing the position away.

Reference genomes and annotations are stored in FASTA (sequences),
GFF or GTF (gene-model annotations), BED (interval features), VCF
(variant calls), and BAM (aligned reads). Each format is a contract
about which fields appear in which order; the engineering discipline
of working in genomics is the discipline of reading the format
specification before reading the data.

\subsection{Pairwise alignment by dynamic programming}

The pairwise-alignment problem asks: given two sequences $a$ and $b$,
what is the maximum-scoring alignment that maps each symbol of one to
a symbol of the other or to a gap, under a fixed substitution score
and gap penalty? Needleman and Wunsch \cite{paper:v6c11-needleman-wunsch-1970}
solved the global version of this problem with dynamic programming in
$O(mn)$ time and $O(mn)$ space, where $m$ and $n$ are the sequence lengths.
The recurrence on the score matrix $M$ is
\[
M[i,j] = \max\!\Bigl\{ M[i-1,j-1] + s(a_i, b_j),\, M[i-1,j] + g,\, M[i,j-1] + g \Bigr\},
\]
with $M[0,0] = 0$ and the first row and column initialised to cumulative
gap penalties: $M[i,0] = i\,g$, $M[0,j] = j\,g$. The score $s(a_i, b_j)$
is a substitution score (positive for matches, negative for mismatches),
and $g$ is a linear gap penalty.

Smith and Waterman \cite{paper:v6c11-smith-waterman-1981} adapted the
recurrence for local alignment by adding a zero floor:
\[
M[i,j] = \max\!\Bigl\{ 0,\, M[i-1,j-1] + s(a_i, b_j),\, M[i-1,j] + g,\, M[i,j-1] + g \Bigr\}.
\]
The local maximum anywhere in the matrix is the best local-alignment
score; traceback starts from that maximum and stops at the first zero.

The recurrence is one of the cleanest illustrations of dynamic
programming in any engineering domain. The state $(i,j)$ stores the
best score for any alignment ending at position $i$ in $a$ and $j$ in
$b$. Three predecessors capture the three ways to extend an alignment
by one column. The optimum is a maximisation over the three. The
correctness rests on the fact that an optimal alignment of the prefixes
must be the prefix of an optimal alignment of the wholes, which is
the discrete-time analogue of Bellman's principle.

\Cref{fig:vol06:ch11:nw-matrix} traces the Needleman-Wunsch recurrence
on the textbook pair GATTACA against GCATGCU with match $+1$, mismatch
$-1$, and gap $-1$. The bottom-right cell holds the optimal global
score; reading the traceback backwards recovers an optimal alignment.

% Figure: Needleman-Wunsch dynamic-programming scoring matrix for the
% pair GATTACA / GCATGCU with match +1, mismatch -1, gap -1.
\begin{figure}[htbp]
\centering
\begin{tikzpicture}[
  cell/.style={draw=engblack!40, minimum width=0.7cm, minimum height=0.6cm,
    font=\small, anchor=center},
  hdr/.style={font=\small\bfseries, anchor=center},
  back/.style={->, thick, draw=engred!70},
]

% Column headers
\node[hdr] at (0.5, 4.0) {};
\node[hdr] at (1.4, 4.0) {$\epsilon$};
\node[hdr] at (2.1, 4.0) {G};
\node[hdr] at (2.8, 4.0) {C};
\node[hdr] at (3.5, 4.0) {A};
\node[hdr] at (4.2, 4.0) {T};
\node[hdr] at (4.9, 4.0) {G};
\node[hdr] at (5.6, 4.0) {C};
\node[hdr] at (6.3, 4.0) {U};

% Row headers
\node[hdr] at (0.5,  3.4) {$\epsilon$};
\node[hdr] at (0.5,  2.8) {G};
\node[hdr] at (0.5,  2.2) {A};
\node[hdr] at (0.5,  1.6) {T};
\node[hdr] at (0.5,  1.0) {T};
\node[hdr] at (0.5,  0.4) {A};
\node[hdr] at (0.5, -0.2) {C};
\node[hdr] at (0.5, -0.8) {A};

% Matrix values
\foreach \j/\v in {0/0, 1/-1, 2/-2, 3/-3, 4/-4, 5/-5, 6/-6, 7/-7} {
  \node[cell, fill=engblue!10] at (1.4 + 0.7*\j, 3.4) {\v};
}
\node[cell, fill=engblue!10] at (1.4,  2.8) {-1};
\node[cell, fill=engblue!10] at (1.4,  2.2) {-2};
\node[cell, fill=engblue!10] at (1.4,  1.6) {-3};
\node[cell, fill=engblue!10] at (1.4,  1.0) {-4};
\node[cell, fill=engblue!10] at (1.4,  0.4) {-5};
\node[cell, fill=engblue!10] at (1.4, -0.2) {-6};
\node[cell, fill=engblue!10] at (1.4, -0.8) {-7};

% Interior cells (selected values shown; computed via recurrence)
\node[cell] at (2.1,  2.8) {1};
\node[cell] at (2.8,  2.8) {0};
\node[cell] at (3.5,  2.8) {-1};
\node[cell] at (4.2,  2.8) {-2};
\node[cell] at (4.9,  2.8) {-1};
\node[cell] at (5.6,  2.8) {-2};
\node[cell] at (6.3,  2.8) {-3};

\node[cell] at (2.1,  2.2) {0};
\node[cell] at (2.8,  2.2) {0};
\node[cell] at (3.5,  2.2) {1};
\node[cell] at (4.2,  2.2) {0};
\node[cell] at (4.9,  2.2) {-1};
\node[cell] at (5.6,  2.2) {-2};
\node[cell] at (6.3,  2.2) {-3};

\node[cell] at (2.1,  1.6) {-1};
\node[cell] at (2.8,  1.6) {-1};
\node[cell] at (3.5,  1.6) {0};
\node[cell] at (4.2,  1.6) {2};
\node[cell] at (4.9,  1.6) {1};
\node[cell] at (5.6,  1.6) {0};
\node[cell] at (6.3,  1.6) {-1};

\node[cell] at (2.1,  1.0) {-2};
\node[cell] at (2.8,  1.0) {-2};
\node[cell] at (3.5,  1.0) {-1};
\node[cell] at (4.2,  1.0) {1};
\node[cell] at (4.9,  1.0) {1};
\node[cell] at (5.6,  1.0) {0};
\node[cell] at (6.3,  1.0) {-1};

\node[cell] at (2.1,  0.4) {-3};
\node[cell] at (2.8,  0.4) {-3};
\node[cell] at (3.5,  0.4) {-1};
\node[cell] at (4.2,  0.4) {0};
\node[cell] at (4.9,  0.4) {0};
\node[cell] at (5.6,  0.4) {0};
\node[cell] at (6.3,  0.4) {-1};

\node[cell] at (2.1, -0.2) {-4};
\node[cell] at (2.8, -0.2) {-2};
\node[cell] at (3.5, -0.2) {-2};
\node[cell] at (4.2, -0.2) {-1};
\node[cell] at (4.9, -0.2) {-1};
\node[cell] at (5.6, -0.2) {1};
\node[cell] at (6.3, -0.2) {0};

\node[cell, fill=englime!25] at (2.1, -0.8) {-5};
\node[cell, fill=englime!25] at (2.8, -0.8) {-3};
\node[cell, fill=englime!25] at (3.5, -0.8) {-1};
\node[cell, fill=englime!25] at (4.2, -0.8) {-2};
\node[cell, fill=englime!25] at (4.9, -0.8) {-2};
\node[cell, fill=englime!25] at (5.6, -0.8) {0};
\node[cell, fill=englime!25] at (6.3, -0.8) {0};

% Optimal score: bottom-right cell highlighted
\node[draw=engred, line width=1.0pt, minimum width=0.7cm,
      minimum height=0.6cm] at (6.3, -0.8) {};

\node[anchor=north west, font=\small, align=left] at (-0.2, -1.6)
  {Recurrence:
   $M[i,j]=\max\{M[i-1,j-1]+s(a_i,b_j),\,M[i-1,j]+g,\,M[i,j-1]+g\}$,\\
   with match $s=+1$, mismatch $s=-1$, gap $g=-1$.
   Optimal global score: $M[m,n]=0$ in this example.};

\end{tikzpicture}
\caption[Needleman-Wunsch DP matrix]{Dynamic-programming
scoring matrix for global alignment of GATTACA against GCATGCU with
match $+1$, mismatch $-1$, and linear gap $-1$. The first row and
column carry cumulative gap penalties (light shaded). The interior
cell at $(i,j)$ takes the maximum of the diagonal predecessor (a
substitution), the cell above (a gap in the second sequence), and
the cell to the left (a gap in the first sequence). The optimal
global score sits in the bottom-right cell; a traceback from that
cell to the top-left recovers an optimal alignment.}
\label{fig:vol06:ch11:nw-matrix}
\end{figure}


\subsection{Affine gap penalties}

The linear-gap model penalises a gap of length $k$ by $k\,g$. In real
biology, gaps come from insertions and deletions whose opening cost is
larger than their extension cost: opening a gap requires a single
mutational event; extending the gap by one residue may cost much less.
The affine-gap model uses two parameters $\alpha$ (gap open) and $\beta$
(gap extend), with total cost $\alpha + (k-1)\,\beta$ for a gap of
length $k\geq 1$. The Gotoh algorithm extends Needleman-Wunsch to
affine gaps by maintaining three score matrices: $M$ for alignments
ending in a match-or-mismatch, $X$ for alignments ending in a gap in
$b$, and $Y$ for alignments ending in a gap in $a$. The recurrence
becomes
\begin{align*}
M[i,j] &= \max\{ M[i-1,j-1], X[i-1,j-1], Y[i-1,j-1] \} + s(a_i, b_j), \\
X[i,j] &= \max\{ M[i-1,j] - \alpha,\, X[i-1,j] - \beta \}, \\
Y[i,j] &= \max\{ M[i,j-1] - \alpha,\, Y[i,j-1] - \beta \}.
\end{align*}
The optimal score sits in $\max\{M[m,n], X[m,n], Y[m,n]\}$. The cost
remains $O(mn)$ in time and space.

\subsection{Substitution matrices}

For protein alignment, the substitution score $s(a_i, b_j)$ is read from
an empirical matrix. PAM matrices (Dayhoff, 1978) were derived from
closely related protein alignments and extrapolated; BLOSUM matrices
\cite{paper:v6c11-henikoff-blosum-1992} were derived from blocks of
conserved regions in less-related proteins. BLOSUM62, the default for
most current BLAST searches, encodes a log-odds score:
\[
s(a, b) = \frac{1}{\lambda} \log \frac{P(a, b)}{P(a) P(b)},
\]
where $P(a, b)$ is the joint frequency of $a$ and $b$ at aligned
positions in conserved blocks and $P(a) P(b)$ is the frequency expected
under independence. The denominator is the null model: the score is
positive when the observed pair appears more often than chance, and
negative when it appears less. The factor $\lambda$ scales the matrix
to integer entries; BLOSUM62 uses $\lambda \approx 0.318$ per nat, so
entries near $\pm 10$ correspond to log-odds ratios of order $e^3$.

\subsection{Heuristic search: BLAST}

The dynamic-programming algorithm is $O(mn)$, which becomes prohibitive
when one of the sequences is a multi-billion-residue database. BLAST
\cite{paper:v6c11-altschul-blast-1990} replaces the exhaustive
recurrence with a seed-and-extend heuristic. The algorithm extracts
all $w$-letter words from the query (with $w=3$ for protein, $w=11$
for DNA), enumerates all words that score above a threshold $T$ when
aligned against the query word, looks up those words in a database
index, and extends each seed hit by ungapped and then gapped local
alignment. The price is a probability of missing significant hits;
the prize is a roughly thousand-fold speedup at typical database
sizes.

The statistical significance of a BLAST hit is given by the
Karlin-Altschul E-value \cite{paper:v6c11-karlin-altschul-1990}:
\[
E = K \, m \, n \, \exp(-\lambda S),
\]
where $S$ is the raw alignment score in score-system units, $m$ is the
query length, $n$ is the total database size, and $K$ and $\lambda$
are Karlin-Altschul parameters that depend only on the scoring system
and the residue composition. The E-value is the expected number of
distinct hits at score $S$ or higher under the null model of random
sequences. A bit score normalises the raw score to a scoring-system-
independent value:
\[
S' = \frac{\lambda S - \ln K}{\ln 2}, \qquad E = \frac{m\,n}{2^{S'}}.
\]
The bit score is portable across BLAST runs; the E-value depends on
both $m$ and $n$ and is read with both numbers in mind.

The E-value formula has an engineering consequence that catches
students of bioinformatics off guard: doubling the database size
doubles the E-value at fixed score, and the same hit that was
significant against a 200~Mb reference becomes routine against a 50~Gb
metagenomic database (\Cref{fig:vol06:ch11:blast-evalue}). The same
sequence has not changed; the inferential weight of finding it by
chance has changed because the search space has grown.

% Figure: BLAST E-value vs database size at fixed query length and
% fixed raw score, showing the linear dependence on m * n.
\begin{figure}[htbp]
\centering
\begin{tikzpicture}
\begin{semilogyaxis}[
  width=11cm, height=7cm,
  xlabel={Database size $n$ (residues, $\times 10^{9}$)},
  ylabel={E-value},
  xmin=0, xmax=10,
  ymin=1e-10, ymax=1e2,
  ymode=log,
  grid=both,
  major grid style={draw=engblack!20},
  minor grid style={draw=engblack!10},
  legend pos=south east,
  legend style={font=\small, draw=engblack!40},
]

% E = K m n exp(-lambda S); fix m = 300, K = 0.1, lambda = 0.3
% Three score curves: S = 60, 80, 100, 120
\addplot[domain=0.01:10, samples=80, color=engblue, thick]
  {0.1 * 300 * x * 1e9 * exp(-0.3 * 60)};
\addlegendentry{$S=60$ bits raw}

\addplot[domain=0.01:10, samples=80, color=englime!80!black, thick]
  {0.1 * 300 * x * 1e9 * exp(-0.3 * 80)};
\addlegendentry{$S=80$ bits raw}

\addplot[domain=0.01:10, samples=80, color=engamber, thick]
  {0.1 * 300 * x * 1e9 * exp(-0.3 * 100)};
\addlegendentry{$S=100$ bits raw}

\addplot[domain=0.01:10, samples=80, color=engred, thick]
  {0.1 * 300 * x * 1e9 * exp(-0.3 * 120)};
\addlegendentry{$S=120$ bits raw}

% Significance threshold
\draw[densely dashed, draw=engblack!60] (axis cs:0,0.01) -- (axis cs:10,0.01);
\node[anchor=west, font=\footnotesize, color=engblack] at (axis cs:0.1,0.03)
  {Common reporting threshold $E=10^{-2}$};

\end{semilogyaxis}
\end{tikzpicture}
\caption[BLAST E-value vs database size]{Karlin-Altschul E-value
$E = K m n \exp(-\lambda S)$ as a function of database size $n$ at
fixed query length $m=300$ residues and Karlin-Altschul parameters
$K=0.1$, $\lambda=0.3$ per bit. Doubling the database doubles the
E-value at fixed score. Each $1/\lambda$ bits of additional raw
score divide the E-value by $e$. The same raw score that is
significant against a 200~Mb genome becomes routine against a 50~Gb
metagenome database; the BLAST result must be re-read every time
the database grows.}
\label{fig:vol06:ch11:blast-evalue}
\end{figure}


Profile-HMM methods (HMMER \cite{paper:v6c11-eddy-hmmer-2011}) replace
the fixed substitution matrix with a position-specific scoring matrix
that captures conserved-versus-variable positions in a protein family,
and replace the linear-gap model with a position-specific gap model.
Profile HMMs are the workhorse of remote-homologue detection, where a
direct BLAST search returns nothing because the pairwise identity is
below the substitution-matrix threshold but a family-level pattern is
preserved.

\begin{archetype}[Information: structure is the signal]
A biological sequence is a string drawn from a small alphabet, but
the strings that biology produces are not uniform over that alphabet.
Real sequences encode coding regions, regulatory regions, structured
RNA, repeats, and lineage-specific signatures. The structure is the
signal: a uniform random string of the same length has Shannon entropy
$\log_2 |\Sigma|$ per symbol (2~bits for DNA, $\approx 4.32$~bits for
protein), but a real coding sequence has entropy on the order of
$1.6$~bits per nucleotide and a homologue's mutual information with
its parent is much larger than the entropy of either alone. The
engineer's task throughout this chapter is to design algorithms that
amplify the structural signal against the alphabet-uniform noise floor.
\end{archetype}

\begin{estimation}
\textbf{Question.}
Estimate the cost, in dollars and in storage, of sequencing and storing
every human genome on Earth by the end of the decade
(current as of 2024).

The human population is approximately $8.0 \times 10^{9}$ persons. A
30-fold-coverage whole-genome sequence at current Illumina pricing
(approximately \$200 per genome amortised including sample prep, current
as of late 2024) gives a per-person cost of order \$200 in pure reagent
terms; clinical-grade processing roughly triples this. At
$\$200 \times 8 \times 10^{9}$ the bare reagent bill is $1.6$~trillion
USD, comparable to the annual healthcare budget of a large economy.

A 30-fold-coverage human genome occupies about 100~GB of FASTQ raw,
which compresses to roughly 30~GB with reference-based compression, and
about 1~MB of meaningful variant calls. Storing the variants for the
population takes $8 \times 10^{9} \times 1\,\text{MB} = 8$~PB. Storing
the raw FASTQ takes $8 \times 10^{9} \times 30\,\text{GB} = 2.4 \times
10^{20}$~bytes, that is, $240$~EB, exceeding by an order of magnitude
the entire 2024 cloud-storage capacity offered by the three largest
public cloud providers combined. The conclusion is operational: the
field moved during the late 2010s from storing reads to storing
variants for routine clinical use, and reverts to the raw reads only
when a reanalysis with a new reference or a new caller is required
\cite{paper:v6c11-stephens-genomic-2015,paper:v6c11-marx-storage-2013}.
\end{estimation}

\section{Sequence-to-structure: prediction, AlphaFold and beyond}\pathtag{standard}

\subsection{The folding problem}

A protein's three-dimensional structure determines its function, and
its structure is determined by its sequence under physiological
conditions. The folding problem asks for the inverse map: given the
amino-acid sequence, predict the atomic coordinates of the native
structure. The problem is in principle a free-energy minimisation over
a configuration space of $10^{2}$ to $10^{4}$ residues, each with a
backbone and a side chain, and the answer is the global minimum of a
free-energy landscape that the cell traverses in milliseconds. The
combinatorial intuition due to Levinthal makes the brute-force search
hopeless: even a 100-residue protein with three rotamers per residue
has $3^{100} \approx 10^{47}$ configurations. The cell does not search;
the cell uses a landscape that funnels efficiently.

\subsection{Physics, statistics, coevolution}

The first generation of structure-prediction methods used physics-based
force fields (CHARMM, AMBER, OPLS) with molecular-dynamics or
Monte-Carlo sampling. Force fields parameterise bond, angle, dihedral,
van der Waals, and electrostatic energies; sampling is the engineering
challenge. The second generation added statistical potentials derived
from the distribution of distances and angles in known structures
(DOPE, Rosetta's terms). The third generation added coevolution
signals: when two residues are in contact in the folded protein,
mutations at one residue tend to be compensated by mutations at the
contacting residue across orthologous sequences, so a multiple-sequence
alignment encodes structural contacts in its correlation pattern.

Direct Coupling Analysis (DCA) and its successors extract pairwise
contact probabilities from MSA covariation by fitting a Potts model on
the alignment and reading the inferred coupling parameters. The
coupling matrix's strongest off-diagonal entries correspond to
residues in contact in the folded structure. By the time of CASP12
(2016), MSA-based contact predictors were producing contact maps
sufficient to fold many proteins with a downstream physics model.

\subsection{AlphaFold and the end-to-end approach}

AlphaFold2 \cite{paper:v6c11-jumper-alphafold-2021} replaced the
two-stage pipeline (contact prediction followed by folding) with an
end-to-end neural-network architecture trained on the PDB. The pipeline
(\Cref{fig:vol06:ch11:alphafold-pipeline}) takes the query sequence,
constructs a multiple-sequence alignment of orthologues by HMM search
(MAFFT or HHsuite \cite{paper:v6c11-katoh-mafft-2013}), searches the
PDB for templates, and feeds both into the Evoformer module. The
Evoformer is a stack of 48 attention blocks that alternates attention
over rows and columns of the MSA representation with triangle updates
of the pair representation. The structure module then converts the
pair representation to a sequence of backbone frames via invariant
point attention, and the side chains are placed by a rotamer head.

% Figure: AlphaFold2 sequence-to-structure pipeline schematic.
\begin{figure}[htbp]
\centering
\begin{tikzpicture}[
  block/.style={draw=engblack, rounded corners=2pt, minimum width=2.6cm,
    minimum height=1.0cm, align=center, font=\small},
  arrow/.style={-{Stealth}, thick, draw=engblack!70},
]

\node[block, fill=engblue!10] (seq) at (0, 0) {Query\\sequence};
\node[block, fill=engblue!15] (msa) at (3.2, 0.7)
  {MSA search\\(MAFFT, HHsuite)};
\node[block, fill=engblue!15] (tpl) at (3.2, -0.7)
  {Template search\\(PDB)};

\node[block, fill=englime!15, minimum width=3.2cm] (evo) at (7.0, 0)
  {Evoformer\\48 blocks, attention\\over MSA + pairs};

\node[block, fill=engamber!15, minimum width=3.2cm] (str) at (11.0, 0)
  {Structure module\\IPA + frame update};

\node[block, fill=engred!10, minimum width=2.6cm] (coord) at (14.5, 0)
  {3D atomic\\coordinates};

% Confidence head
\node[block, fill=enggray!15, minimum width=2.6cm] (plddt) at (11.0, -2.0)
  {pLDDT head\\(per-residue\\confidence)};

\draw[arrow] (seq) -- (msa);
\draw[arrow] (seq) -- (tpl);
\draw[arrow] (msa) -- (evo);
\draw[arrow] (tpl) -- (evo);
\draw[arrow] (evo) -- (str);
\draw[arrow] (str) -- (coord);
\draw[arrow] (str.south) -- (plddt.north);
\draw[arrow, draw=engblack!40] (evo.south) |- (plddt.west);

% Recycling loop
\draw[arrow, draw=engred!60, dashed]
  (str.north) to[bend left=30] node[midway, above, font=\footnotesize]
  {recycle ($\times 3$ typical)} (evo.north);

\end{tikzpicture}
\caption[AlphaFold pipeline schematic]{AlphaFold2 pipeline from
sequence to predicted structure. A multiple-sequence alignment
captures coevolutionary information across orthologues; templates
from the PDB provide a starting geometry where homologues exist.
The Evoformer alternates row-attention, column-attention, and
triangle updates between the MSA stack and the pair stack. The
structure module converts the pair representation to a sequence of
backbone frames via invariant point attention. The pLDDT head emits
a per-residue confidence score that is the calibrated check against
the coordinates. The full forward pass is recycled three times
during inference; the cost is approximately linear in target
length for sequences below 2~kAA.}
\label{fig:vol06:ch11:alphafold-pipeline}
\end{figure}


Training combined the PDB ($\sim 175\text{k}$ structures by 2021) with
self-distillation: AlphaFold predicted structures for sequences without
deposited structures, filtered by predicted confidence, and added the
high-confidence predictions back to the training set. The
median backbone accuracy at CASP14 (2020) was $1.0$~A RMSD on
template-free targets, comparable to experimental X-ray crystallography
for the easier targets in the set
\cite[][\ p.~586]{paper:v6c11-jumper-alphafold-2021}.

The per-residue confidence score pLDDT (predicted local distance
difference test) is the calibrated check on the coordinates. Regions
with pLDDT below $50$ are predicted to be disordered or to have low
local accuracy; regions with pLDDT above $90$ are predicted to be at
near-experimental accuracy. Reading an AlphaFold structure without
reading pLDDT is reading a prediction without its uncertainty.

\subsection{ESMFold and the language-model lineage}

A parallel line of work used protein language models (ESM-1, ESM-2)
\cite{paper:v6c11-lin-esm2-2023} trained by masked-language modelling on
hundreds of millions of unaligned protein sequences. ESMFold reads
structure from a single sequence using only the language model's
embeddings, with no MSA. The accuracy is below AlphaFold for targets
with deep MSAs but is competitive for targets without orthologues
(orphans, designed proteins, viral proteins from undersampled lineages),
and is roughly two orders of magnitude faster. The two architectures
embody two information sources for the folding problem: evolution
across orthologues (AlphaFold's MSA), and the implicit grammar of
folded proteins (ESMFold's language model).

\subsection{What remains hard}

Several classes of structure prediction remain difficult as of 2024.
Multi-state proteins (kinases, GPCRs, transporters) fold into many
distinct conformations; AlphaFold typically returns the dominant
conformation in the training distribution, not the one relevant to the
user's question. Intrinsically disordered regions (which make up
$\approx 30\%$ of human protein residues) do not have a unique
structure to predict; pLDDT is low in those regions, which is the
correct answer, but the calling convention "structure of protein X"
becomes ill-defined. Membrane proteins were under-represented in the
PDB training corpus, and the accuracy on transmembrane regions remains
below the accuracy on soluble proteins. Protein-ligand complexes,
metal-binding sites, and post-translational modifications are not
modelled in the base architecture; specialised heads (AlphaFill,
AlphaFold-Multimer, RoseTTAFold All-Atom) extend coverage but at
correspondingly degraded accuracy. The reader takes a pLDDT and the
class of protein into account before quoting an AlphaFold structure
as a fact.

\section{Genomics and variant calling}\pathtag{standard}

\subsection{Reference genomes and short-read alignment}

Modern whole-genome sequencing produces hundreds of millions of short
reads (75-300~bp typical Illumina, $10$-$100$~kb typical PacBio or
Oxford Nanopore) per sample. The reads are aligned to a reference
genome to identify the position from which each read originated. Two
algorithmic ideas dominate short-read alignment: hash-based seeding
(Bowtie-1 era) and FM-index seeding using the Burrows-Wheeler
transform (BWA-MEM \cite{paper:v6c11-li-bwa-2009}, Bowtie 2
\cite{paper:v6c11-langmead-bowtie-2012}).

The FM-index stores the BWT of the reference along with a sparse rank
data structure that supports exact substring search in time linear in
the substring length and independent of the reference length. The
total memory for a human-genome FM-index is approximately $4$~GB, which
fits in memory on a workstation. Once seeds are placed by exact
search, gapped extension uses a banded dynamic-programming alignment
on the seed neighbourhood. The output is a SAM (sequence alignment
map) record per read carrying position, CIGAR string (the alignment
operation sequence), mapping quality, and per-base quality.

\subsection{Variant calling}

A variant call is a position in the reference where the sample
sequence differs from the reference. The simplest variants are
single-nucleotide variants (SNVs) and short insertions or deletions
(indels, up to $\sim 50$~bp). Larger structural variants (inversions,
translocations, copy-number variants) require specialised callers and
are not treated here.

The canonical SNV-and-indel pipeline (GATK
\cite{paper:v6c11-mckenna-gatk-2010}) is a cascade of statistical
steps. The HaplotypeCaller assembles short de Bruijn graphs from reads
in active regions, identifies candidate haplotypes, realigns reads
against the haplotypes, and calls variants from the realigned pileup.
The output is raw variant calls with per-call statistics: QUAL (the
log-odds of the variant being real), DP (total read depth), AF
(alternate-allele fraction), FS (Fisher's exact strand-bias score),
MQ (root-mean-square mapping quality), and others.

Filtering converts raw calls to confident calls
(\Cref{fig:vol06:ch11:variant-pipeline}). Hard filters apply
thresholds: $\mathrm{QUAL} \geq 30$, $\mathrm{DP} \in [10, 200]$ for a
diploid germline call, $\mathrm{FS} \leq 60$, $\mathrm{MQ} \geq 40$.
Variant Quality Score Recalibration (VQSR) fits a Gaussian mixture to
high-confidence training variants and scores new calls by their
posterior probability of belonging to the high-confidence cluster;
this is the recommended approach for population-scale projects.
DeepVariant \cite{paper:v6c11-poplin-deepvariant-2018} reframes the
problem as image classification on pileup images and trains a
convolutional network to call variants directly, achieving
state-of-the-art accuracy on the Genome in a Bottle benchmark.

% Figure: Variant-calling pipeline schematic from raw reads to filtered
% variant set with quality, depth, and strand-bias filters.
\begin{figure}[htbp]
\centering
\begin{tikzpicture}[
  block/.style={draw=engblack, rounded corners=2pt, minimum width=3.4cm,
    minimum height=0.9cm, align=center, font=\small},
  filt/.style={draw=engred!60, rounded corners=2pt, minimum width=3.4cm,
    minimum height=0.9cm, align=center, font=\small, fill=engred!8},
  arrow/.style={-{Stealth}, thick, draw=engblack!70},
  reject/.style={-{Stealth}, thick, draw=engred!60, dashed},
]

\node[block, fill=engblue!10] (raw) at (0, 5)
  {Raw reads\\(FASTQ)};
\node[block, fill=engblue!10] (qc) at (0, 3.6)
  {Quality trimming\\(Phred $\geq 20$)};
\node[block, fill=engblue!10] (align) at (0, 2.2)
  {Alignment\\(BWA-MEM, Bowtie2)};
\node[block, fill=engblue!10] (dedup) at (0, 0.8)
  {Duplicate removal\\(MarkDuplicates)};
\node[block, fill=engblue!10] (call) at (0, -0.6)
  {Raw variant calls\\(GATK, DeepVariant)};

\node[filt] (qual) at (5.0, 0.8) {Quality filter\\$\mathrm{QUAL} \geq 30$};
\node[filt] (depth) at (5.0, -0.6) {Depth filter\\$\mathrm{DP}\in[10,200]$};
\node[filt] (sb) at (5.0, -2.0) {Strand-bias filter\\$\mathrm{FS} \leq 60$};
\node[filt] (mq) at (5.0, -3.4) {Mapping-quality filter\\$\mathrm{MQ}\geq 40$};

\node[block, fill=englime!20, minimum width=3.4cm, minimum height=1.0cm]
  (final) at (10.0, -1.3)
  {Filtered variants\\(VCF, PASS only)};

\node[block, fill=enggray!10, font=\footnotesize] (rej) at (10.0, -3.6)
  {Rejected variants\\(retained, flagged)};

\draw[arrow] (raw) -- (qc);
\draw[arrow] (qc) -- (align);
\draw[arrow] (align) -- (dedup);
\draw[arrow] (dedup) -- (call);
\draw[arrow] (call) -- (qual);
\draw[arrow] (qual) -- (depth);
\draw[arrow] (depth) -- (sb);
\draw[arrow] (sb) -- (mq);
\draw[arrow] (mq.east) -| (final.south);

\draw[reject] (qual.east) -- ++(1.5, 0) |- (rej.west);
\draw[reject] (depth.east) -- ++(1.5, 0) |- (rej.west);
\draw[reject] (sb.east) -- ++(1.5, 0) |- (rej.west);
\draw[reject] (mq.east) -- ++(1.5, 0) |- (rej.west);

\end{tikzpicture}
\caption[Variant-calling pipeline]{Variant-calling pipeline from
raw sequencing reads to a filtered set of confident variant calls.
The pipeline is a cascade of filters, each justified by a known
failure mode: low Phred scores reflect base-call uncertainty;
deduplication removes the over-confidence introduced by PCR
duplicates; the quality filter rejects calls whose variant-vs-error
likelihood ratio is too small; the depth filter rejects calls
supported by too few reads (low confidence) or by too many reads
(repeat or copy-number artefact); the strand-bias filter rejects
calls supported predominantly by one strand (read-orientation
artefact); the mapping-quality filter rejects calls in regions of
ambiguous alignment. Rejected variants are retained in the VCF with
their FILTER field set, not deleted, so downstream tools can
reconsider them.}
\label{fig:vol06:ch11:variant-pipeline}
\end{figure}


The filter cascade is not optional. A modern human whole-genome
sequence yields approximately $4$ million raw single-nucleotide
variants; the filtered set is approximately $3.5$ million, with
discordance between callers running at $5$ to $10\%$ on the more
difficult variant classes. The number that propagates to the clinical
report (the medically actionable variants) is typically below $200$.
Each filter step in the pipeline corresponds to a known failure mode
that, left in, accounts for a fraction of the false positive rate.

\subsection{Variant call format}

Variants are stored in VCF \cite{paper:v6c11-danecek-vcf-2011}. Each
data line carries CHROM, POS, ID, REF, ALT, QUAL, FILTER, INFO, plus
per-sample genotype columns. The FILTER column carries either PASS or
a semicolon-separated list of filter names that the record failed;
rejected variants are retained in the file with their failure tags,
not deleted, so a downstream tool can reconsider them if a different
threshold applies. The INFO column carries semicolon-separated
key=value pairs with the per-call statistics. The discipline of
working in VCF is the discipline of reading the header lines, which
declare what each INFO key means and how the file was produced.

\subsection{Population-scale variant interpretation}

A variant's biological significance is rarely apparent from the
variant itself. The variant is matched against population databases
(gnomAD, 1000 Genomes), against clinical-significance databases
(ClinVar), against gene-annotation databases (Ensembl
\cite{paper:v6c11-ensembl-2023}, RefSeq), and against the
disease-specific cohort, before a clinical interpretation is issued.
The same variant can be benign in a population with $5\%$ minor-allele
frequency and pathogenic in a population where it is private and
co-segregates with disease. The interpretation is a small calculation
on top of a large lookup; the lookup tables are the asset.

\section{Transcriptomics and single-cell RNA-seq}\pathtag{standard}

\subsection{Bulk RNA-seq}

The transcriptome is the set of all RNA molecules present in a sample.
Bulk RNA-seq measures the abundance of each transcript across a
population of cells. The workflow is: extract RNA, deplete ribosomal
RNA or capture polyadenylated transcripts, fragment, reverse-transcribe
to cDNA, attach sequencing adapters, sequence to a target depth (10-30
million reads is typical for differential expression), align reads to
the transcriptome or genome, count reads per transcript or per gene,
normalise across samples, and test for differential expression.

The output of a bulk RNA-seq experiment is a counts matrix: rows are
features (genes or transcripts), columns are samples, entries are read
counts. Counts must be normalised before comparison. Library-size
normalisation divides by total counts per sample; TMM normalisation
(edgeR) and median-of-ratios normalisation (DESeq2) fit a per-sample
size factor under the assumption that most genes are not differentially
expressed; quantile normalisation forces the empirical distribution of
counts to be identical across samples (used when the assumption that
most genes are not differential is questionable).

Differential expression testing uses negative-binomial generalised
linear models (edgeR, DESeq2). The negative binomial captures both the
Poisson sampling component (the read-count process) and the
biological over-dispersion (variance greater than mean across
biological replicates). The dispersion parameter is estimated per
gene with shrinkage toward a pooled estimate, which is the canonical
empirical-Bayes pattern: borrow strength across genes to stabilise the
per-gene variance estimate.

\subsection{Single-cell RNA-seq}

Single-cell RNA-seq (scRNA-seq) measures the transcriptome of each of
many individual cells separately. The workflow uses microfluidic
droplets (10x Genomics Chromium, inDrops, Drop-seq), microwells, or
plate-based methods to capture each cell into its own reaction volume,
attach a cell-specific barcode and a unique molecular identifier
(UMI), and pool the reactions for sequencing. The output is a counts
matrix with cells as columns instead of samples; a typical 10x
Chromium run yields 5000 to 10000 cells per channel.

The data are sparse: a typical cell has 1000 to 5000 captured
transcripts out of 200,000 or more present in the cell, so the counts
matrix is more than $90\%$ zero. Many of those zeros are "missing"
(the transcript was present but not captured); some are "true"
(the transcript was not expressed). The distinction matters for
downstream analysis: differential expression tests that assume zero
means absence will mis-call genes whose only failure was capture
inefficiency. Modern scRNA-seq pipelines model the zeros explicitly
(ZINB-WaVE) or use methods that are robust to zero-inflation (MAST,
Wilcoxon rank-sum).

\subsection{Workflow: QC, normalisation, clustering, integration}

The canonical scRNA-seq workflow \cite{paper:v6c11-stuart-scrnaseq-2019}
proceeds in stages. Quality control removes cells with too few or too
many counted transcripts (likely empty droplets or doublets), removes
cells with high mitochondrial-read fraction (likely stressed or dying
cells), and removes ambient-RNA contamination by SoupX or DecontX.
Normalisation by SCT (Seurat) or scran size-factors corrects for cell-
to-cell capture variation. Highly variable genes are selected to
reduce dimensionality. PCA projects to 30-50 components. A k-nearest-
neighbour graph is constructed in PCA space; the graph is clustered by
Leiden or Louvain. The cluster identities are assigned to cell types
by marker-gene expression. UMAP or t-SNE produces a two-dimensional
visualisation for inspection.

% Figure: Sketch of a scRNA-seq UMAP showing well-separated and
% poorly-separated cell-type clusters.
\begin{figure}[htbp]
\centering
\begin{tikzpicture}
\begin{axis}[
  width=11cm, height=8cm,
  xlabel={UMAP-1},
  ylabel={UMAP-2},
  xmin=-7, xmax=10,
  ymin=-7, ymax=7,
  grid=both,
  major grid style={draw=engblack!15},
  minor grid style={draw=engblack!8},
  ticks=none,
  legend pos=outer north east,
  legend style={font=\footnotesize, draw=engblack!40},
]

% Cluster 1: T cells
\addplot[only marks, mark=*, mark size=0.6pt, color=engblue]
  table {
  x y
  -4.0 4.0
  -3.7 3.5
  -4.2 4.3
  -3.9 4.1
  -3.5 3.7
  -4.4 4.0
  -3.6 4.2
  -4.0 3.6
  -3.8 4.4
  -4.3 3.8
  -3.5 4.1
  -4.1 3.9
  -3.7 4.5
  -4.0 4.2
  -3.9 3.5
  -4.2 4.0
  -3.6 3.8
  -4.1 4.3
  -3.8 3.9
  -4.0 3.7
  };
\addlegendentry{CD4 T (well-separated)}

% Cluster 2: B cells
\addplot[only marks, mark=*, mark size=0.6pt, color=engamber]
  table {
  x y
  5.5 3.5
  5.8 3.0
  5.2 3.8
  5.7 3.4
  5.4 3.2
  5.9 3.6
  5.3 3.7
  5.6 3.1
  5.8 3.5
  5.2 3.3
  5.5 3.9
  5.7 3.7
  5.4 3.4
  5.9 3.2
  5.3 3.6
  5.6 3.8
  };
\addlegendentry{B cells (well-separated)}

% Cluster 3: Monocytes, with sub-structure
\addplot[only marks, mark=*, mark size=0.6pt, color=englime!80!black]
  table {
  x y
  0.0 -3.0
  0.3 -3.3
  -0.2 -2.7
  0.1 -3.1
  0.4 -3.4
  -0.1 -2.9
  0.2 -3.2
  0.5 -3.0
  -0.3 -2.8
  0.0 -3.4
  0.3 -2.8
  -0.2 -3.1
  0.1 -2.9
  0.4 -3.2
  -0.1 -3.4
  0.6 -3.0
  };
\addlegendentry{Monocytes (one cluster)};

% Cluster 4: CD8 T, overlapping CD4 T (illustrates poor separation)
\addplot[only marks, mark=*, mark size=0.6pt, color=engred]
  table {
  x y
  -2.5 3.5
  -2.8 3.0
  -2.2 3.8
  -2.7 3.4
  -2.4 3.2
  -2.9 3.6
  -2.3 3.7
  -2.6 3.1
  -2.8 3.5
  -2.2 3.3
  -3.0 3.9
  -2.5 3.7
  -2.4 3.4
  -2.9 3.2
  -2.3 3.6
  -2.6 3.8
  -2.7 4.0
  };
\addlegendentry{CD8 T (overlaps CD4)};

% Doublets: scattered, not cluster
\addplot[only marks, mark=x, mark size=2pt, color=engblack!60]
  table {
  x y
  1.5 1.0
  -1.0 0.5
  2.5 -0.5
  -2.0 -1.5
  3.5 1.5
  };
\addlegendentry{Putative doublets};

% Annotation labels
\node[anchor=west, font=\footnotesize, align=left] at (axis cs:-7, 5.0)
  {Cluster well-separated\\$\Rightarrow$ marker enrichment\\reliable.};
\node[anchor=west, font=\footnotesize, color=engred!80!black, align=left]
  at (axis cs:-4, 1.0)
  {CD4 and CD8\\overlap:\\marker genes\\required for\\subtype call.};

\end{axis}
\end{tikzpicture}
\caption[scRNA-seq UMAP cartoon]{Sketch of a UMAP embedding of a
small scRNA-seq cohort. Well-separated clusters (B cells, CD4 T
cells, monocytes) correspond to populations whose marker-gene
expression differs strongly. CD4 and CD8 T cells overlap in the
two-dimensional projection because their transcriptional signatures
differ in a small number of canonical markers (CD4, CD8A, CD8B)
against a large shared T-cell background; the cell-type call rests
on those markers rather than on the embedding coordinates. Putative
doublets (cells where two distinct cell-type signatures appear in
one droplet) scatter outside the clusters and are flagged by
preprocessing tools rather than by UMAP. The UMAP coordinates are
a visualisation aid, not a metric; cluster assignment is performed
on the upstream Leiden or graph-based clustering of the high-
dimensional neighbour graph.}
\label{fig:vol06:ch11:scrna-umap}
\end{figure}


\Cref{fig:vol06:ch11:scrna-umap} shows the failure mode common to
UMAP interpretation. Well-separated clusters in the embedding
correspond to populations whose marker genes differ strongly. Clusters
that overlap in UMAP may differ in a small number of canonical
markers against a large shared transcriptional background; the
cell-type call rests on the markers, not on the embedding coordinates.
Reading a UMAP plot as a metric is the canonical first error of a
new scRNA-seq analyst. The two-dimensional UMAP preserves local
neighbourhoods at the expense of global geometry; distances in the
plot do not faithfully represent distances in transcript space.

\subsection{Integration across batches and conditions}

scRNA-seq experiments are routinely conducted across multiple batches,
laboratories, time points, or conditions. Integration aligns the
common cell-type axis across batches while preserving the
between-condition differences. Anchor-based methods (Seurat v3,
SCMAP) find mutual nearest neighbours across batches in PCA space and
use them to estimate batch-correction vectors. Variational autoencoder
methods (scVI, scANVI) model the joint distribution and learn a latent
representation in which batch is regressed out. The integration step
is where the disease-vs-batch problem of the failure section appears:
if a condition is collinear with batch, no integration can recover the
condition signal honestly.

\section{Proteomics, metabolomics, lipidomics, fluxomics}\pathtag{standard}

\subsection{The omics stack}

\Cref{fig:vol06:ch11:omics-stack} sketches the omics stack as a chain
of measurement layers from genome (static, $\sim 3 \times 10^9$ bp in
human) to phenotype (the system-level outcome). Each layer is
measured by its own technology with its own artefact class. A study
that uses only one layer measures the projection of a many-dimensional
state onto a single axis; multi-omics integration estimates the joint
state from partial views with very different noise statistics.

% Figure: The omics stack: genome, transcriptome, proteome, metabolome,
% lipidome, fluxome, showing what each layer measures and at what
% scale.
\begin{figure}[htbp]
\centering
\begin{tikzpicture}[
  layer/.style={draw=engblack, rounded corners=2pt, minimum width=4.0cm,
    minimum height=0.9cm, align=center, font=\small},
  ann/.style={anchor=west, font=\footnotesize, color=enggray},
  arrow/.style={->, thick, draw=engblack!60},
]

% Vertical stack of layers (top: genome; bottom: phenotype)
\node[layer, fill=engblue!10]   (g) at (0, 5.0) {Genome};
\node[layer, fill=engblue!20]   (t) at (0, 3.7) {Transcriptome};
\node[layer, fill=englime!20]   (p) at (0, 2.4) {Proteome};
\node[layer, fill=engamber!20]  (m) at (0, 1.1) {Metabolome / Lipidome};
\node[layer, fill=engred!15]    (f) at (0, -0.2) {Fluxome};
\node[layer, fill=enggray!20]   (ph) at (0, -1.5) {Phenotype};

% Right-side annotations: what is measured and how
\node[ann, align=left] at (3.0, 5.0)
  {$\sim 3 \times 10^{9}$~bp; static.\\
   Measured by short-read sequencing,\\
   stored as FASTA + VCF.};
\node[ann, align=left] at (3.0, 3.7)
  {$\sim 2 \times 10^{4}$~genes, $10^{5}$~transcript isoforms.\\
   Dynamic, condition-specific.\\
   Measured by RNA-seq, scRNA-seq.};
\node[ann, align=left] at (3.0, 2.4)
  {$\sim 2 \times 10^{4}$~proteins, $10^{6}$~proteoforms.\\
   Post-translational state.\\
   Measured by mass spectrometry.};
\node[ann, align=left] at (3.0, 1.1)
  {$\sim 10^{4}$~metabolites, $10^{3}$~lipid species.\\
   Dynamic, fast-turnover.\\
   Measured by LC-MS, GC-MS, NMR.};
\node[ann, align=left] at (3.0, -0.2)
  {$\sim 10^{3}$~reactions in central metabolism.\\
   Inferred, not measured directly.\\
   FBA, $^{13}$C tracing, kinetic models.};
\node[ann, align=left] at (3.0, -1.5)
  {Growth, survival, disease, fitness.\\
   The reason the rest of the stack exists.};

% Vertical arrows
\draw[arrow] (g.south) -- (t.north);
\draw[arrow] (t.south) -- (p.north);
\draw[arrow] (p.south) -- (m.north);
\draw[arrow] (m.south) -- (f.north);
\draw[arrow] (f.south) -- (ph.north);

% Left-side process labels
\node[font=\footnotesize\itshape, color=enggray, anchor=east, align=right]
  at (-2.2, 4.35) {transcription\\+ regulation};
\node[font=\footnotesize\itshape, color=enggray, anchor=east, align=right]
  at (-2.2, 3.05) {translation\\+ PTM};
\node[font=\footnotesize\itshape, color=enggray, anchor=east, align=right]
  at (-2.2, 1.75) {enzymatic\\activity};
\node[font=\footnotesize\itshape, color=enggray, anchor=east, align=right]
  at (-2.2, 0.45) {flux distribution\\through network};
\node[font=\footnotesize\itshape, color=enggray, anchor=east, align=right]
  at (-2.2, -0.85) {organism-level\\integration};

\end{tikzpicture}
\caption[The omics stack]{The omics stack reads bottom-up as the
chain of causation by which genome becomes phenotype, and top-down
as the chain of measurement by which engineers reconstruct what
the cell is doing. Each layer has its own measurement technology,
its own scale, its own dynamic range, and its own characteristic
artefact class. A study that measures only one layer measures the
projection of a many-dimensional state onto a single axis; the
multi-omics integration problem is the engineering problem of
estimating the joint state from partial views with very different
noise statistics. Counts and figures are order-of-magnitude human
estimates current as of 2024.}
\label{fig:vol06:ch11:omics-stack}
\end{figure}


\subsection{Proteomics by mass spectrometry}

Mass-spectrometry-based proteomics \cite{paper:v6c11-aebersold-proteomics-2003}
measures the abundance of proteins in a sample by digesting proteins to
peptides (typically with trypsin), separating peptides by
liquid chromatography, ionising them, and measuring their mass-to-charge
ratio in a mass spectrometer. Modern instruments perform tandem MS:
the precursor peptide ion is selected by its m/z, fragmented by
collision-induced dissociation, and the fragment-ion spectrum is
recorded. The fragment-ion pattern is a sequence-specific fingerprint
of the peptide.

Peptide identification by database search (SEQUEST, MaxQuant, MSFragger)
compares each observed spectrum against the theoretical spectra of all
peptides predicted from a digest of a reference proteome. The score is
typically a hypergeometric or cross-correlation match. False discovery
rate is controlled by target-decoy: searching the database concatenated
with a shuffled or reversed copy and counting decoy matches at each
score threshold to estimate the FDR among target matches above that
threshold. The conventional reporting threshold is $1\%$ peptide FDR
followed by $1\%$ protein FDR.

Quantification approaches divide into label-free (the area under the
chromatographic peak, normalised across runs), label-based (isobaric
tags such as iTRAQ or TMT that allow multiplexed comparison in a
single MS run), and SILAC (stable-isotope labelling by amino acids in
cell culture, which compares two or three differentially labelled
samples in the same instrument run). Each method has its own
linearity range and its own batch-effect class.

Post-translational modifications (PTMs) are read by mass shifts at
the modification site: phosphorylation adds $+80$~Da, methylation
$+14$~Da, ubiquitination leaves a $+114$~Da GG remnant after digestion.
The combinatorial explosion of PTMs ($\sim 10^6$ proteoforms in a
typical human proteome) is the upstream cause of the gap between
proteome size ($\sim 2 \times 10^{4}$ canonical proteins) and
proteoform space.

\subsection{Metabolomics and lipidomics}

The metabolome is the set of small molecules ($<1500$~Da) present in a
sample. The lipidome is the lipid-specific subset. Both are measured
by mass spectrometry (LC-MS, GC-MS) or by NMR. Targeted assays
quantify a known list of metabolites with internal standards; untargeted
assays scan the whole observable mass range and identify peaks
post-hoc. The Human Metabolome Database \cite{paper:v6c11-wishart-hmdb-2022}
catalogues approximately $250,000$ entries (current as of 2022),
of which on the order of $5{,}000$ are routinely detected in plasma.

The lipidomics taxonomy distinguishes glycerolipids, glycerophospholipids,
sphingolipids, sterol lipids, prenol lipids, saccharolipids, polyketides,
and fatty acyls \cite{paper:v6c11-han-lipidomics-2016}. Lipid identification
requires both the precursor m/z and the head-group fragment-ion pattern;
chain-length and double-bond position ambiguity remains the dominant
identification uncertainty in shotgun lipidomics.

\subsection{Fluxomics and flux balance analysis}

The fluxome is the set of metabolic-reaction rates in a cell at steady
state. Direct measurement uses isotope tracing ($^{13}$C-labelled
glucose or amino acids, with the labelling pattern at downstream
metabolites used to infer the flux through each pathway). Indirect
estimation uses flux balance analysis (FBA \cite{paper:v6c11-orth-cobra-2010})
on a constraint-based model of the metabolic network.

The FBA formulation is a linear program. The stoichiometric matrix
$S$ has rows for metabolites and columns for reactions. At steady state,
the flux vector $v$ satisfies $S v = 0$. Reaction reversibility imposes
lower and upper bounds, $v^{\mathrm{lb}} \leq v \leq v^{\mathrm{ub}}$.
The objective is a linear function of the fluxes, typically the
biomass-production rate (a synthetic reaction whose stoichiometry
encodes the cellular composition needed for one unit of growth). The
LP is
\[
\max_{v}\; c^{\top} v \quad \text{subject to} \quad S v = 0, \quad
v^{\mathrm{lb}} \leq v \leq v^{\mathrm{ub}}.
\]
The solution is the flux distribution that maximises growth under the
network's stoichiometric constraints. FBA is a single LP; it scales to
genome-scale models with $\sim 10^{3}$ to $10^{4}$ reactions in
seconds.

The interpretive caveats are large. FBA assumes steady state; many
biologically interesting transitions are not at steady state. FBA
assumes the cell maximises growth; many cells in many contexts do not.
FBA returns a flux distribution that is optimal under the assumed
objective, not necessarily the actual flux distribution; flux
variability analysis quantifies the range of fluxes consistent with
the optimal objective. Thermodynamic feasibility adds further
constraints \cite{paper:v6c11-noor-fluxomics-2014}: each reaction must
have $\Delta G < 0$ in the direction of the flux. The FBA solution
becomes credible when the constraints, the objective, and the steady-
state assumption are each defended for the biological context.

\begin{estimation}
\textbf{Question.}
Estimate the compute cost of folding every protein in the human
proteome with AlphaFold2 on commercial cloud GPUs (current as of 2024).

The human proteome contains approximately $2 \times 10^{4}$ canonical
proteins with average length $\sim 400$ residues. AlphaFold2 inference
scales roughly as $O(L^{2})$ in attention; on an A100 GPU, a 400-residue
target takes approximately $2$ minutes including MSA construction. A
shorter target may take $30$ seconds, a longer one (1000+ residues)
may take an hour or more.

Total compute: $2 \times 10^{4}$ proteins $\times$ $2$ minutes
$= 4 \times 10^{4}$~min $\approx 670$~GPU-hours on A100. At
commercial cloud spot pricing of approximately \$1.50 per A100-hour
(current as of 2024), the total cost is approximately \$1{,}000 for
the canonical proteome, plus the MSA-search cost (which adds memory
and CPU time but is dominated by the GPU bill for short to medium
proteins). DeepMind and the EBI \cite{web:v6c11-alphafold-db-2024}
published a free database that contains approximately $214$ million
predicted structures spanning UniProt as of 2024; the cost of the
underlying compute is on the order of hundreds of thousands of
GPU-hours. The order-of-magnitude conclusion is that proteome-scale
folding is now within the reach of small laboratories at a per-proteome
cost of $\sim 10^{3}$ USD; folding all of UniProt at $\sim 2.5 \times
10^{8}$ entries requires substantially more compute and remains a
shared-infrastructure project.
\end{estimation}

\section{Phylogenetics and evolutionary engineering}\pathtag{standard}

\subsection{The phylogeny problem}

Phylogenetics asks for the tree that best describes the evolutionary
relationships among a set of taxa, given their molecular sequences.
The tree has the taxa as leaves and internal nodes that correspond to
hypothetical ancestral lineages; the tree's branch lengths represent
amounts of evolutionary change between nodes. Three method families
dominate: distance methods (UPGMA, neighbour-joining
\cite{paper:v6c11-saitou-nj-1987}), parsimony methods (find the tree
minimising the number of inferred substitutions), and likelihood or
Bayesian methods (find the tree maximising the probability of the
data under an explicit substitution model
\cite{paper:v6c11-felsenstein-phylogeny-1981}).

\subsection{Distance methods}

Distance methods reduce the multiple-sequence alignment to a pairwise-
distance matrix and reconstruct the tree from the distances.
UPGMA (unweighted pair group method with arithmetic mean) iteratively
joins the closest pair of clusters, weighting each cluster's
contribution to the merged distance by its size. The resulting tree is
ultrametric: every leaf is at the same distance from the root, which is
consistent with a strict molecular clock and is rarely a good model for
real lineages. Neighbour-joining drops the ultrametric assumption and
produces an additive tree that minimises the total branch length;
neighbour-joining is provably consistent under the assumption that the
true distances are additive and that the distance estimates are
asymptotically unbiased.

The substitution-distance estimate corrects raw mismatch fractions
for multiple substitutions at the same position. Under the Jukes-Cantor
model (uniform substitution rate), the corrected distance from a
mismatch fraction $p$ is
\[
d = -\frac{3}{4} \ln\!\left(1 - \frac{4}{3}p\right).
\]
The Kimura two-parameter model distinguishes transitions from
transversions; the HKY and GTR models add unequal base frequencies and
unequal substitution rates.

\Cref{fig:vol06:ch11:phylo-tree} shows a five-taxon maximum-likelihood
phylogeny from a 1.5~kb rRNA alignment, with branch lengths in
substitutions per site and bootstrap support values on the internal
edges. The bootstrap is a resampling procedure: the alignment columns
are resampled with replacement to produce a new alignment of the same
length, a tree is inferred from each resample, and the fraction of
resamples that recover each internal partition is the bootstrap
support. Values above 80 are commonly taken as strong evidence for the
partition; values below 70 are not.

% Figure: Phylogenetic tree with branch lengths and bootstrap support
% values for a small five-taxon 16S example.
\begin{figure}[htbp]
\centering
\begin{tikzpicture}[
  taxon/.style={anchor=west, font=\small},
  ilabel/.style={font=\footnotesize\itshape, color=enggray},
  brlen/.style={font=\footnotesize, color=engblue!80!black},
  bsval/.style={font=\footnotesize, color=engred!80!black},
  edge/.style={thick, draw=engblack!80},
]

% Vertical positions for the five tips
% Tip order (top to bottom): Aspergillus, Saccharomyces, E.coli, B.subtilis, Halobacterium

% Tips
\node[taxon] (a) at (8, 4.0)  {\textit{Aspergillus niger}};
\node[taxon] (s) at (8, 3.2)  {\textit{Saccharomyces cerevisiae}};
\node[taxon] (e) at (8, 2.0)  {\textit{Escherichia coli}};
\node[taxon] (b) at (8, 1.2)  {\textit{Bacillus subtilis}};
\node[taxon] (h) at (8, 0.0)  {\textit{Halobacterium salinarum}};

% Internal nodes
\coordinate (n_as) at (6.5, 3.6);   % a - s
\coordinate (n_eb) at (6.0, 1.6);   % e - b
\coordinate (n_euk) at (4.5, 3.6);  % eukaryote clade
\coordinate (n_bac) at (4.5, 1.6);  % bacteria clade
\coordinate (n_root) at (2.0, 2.6); % deep root with archaea

% Edges to tips
\draw[edge] (n_as) -- node[brlen, above] {0.05} (a.west);
\draw[edge] (n_as) -- node[brlen, below] {0.08} (s.west);
\draw[edge] (n_eb) -- node[brlen, above] {0.12} (e.west);
\draw[edge] (n_eb) -- node[brlen, below] {0.11} (b.west);
\draw[edge] (n_root) -- node[brlen, below, sloped] {0.28} (h.west);

% Internal edges
\draw[edge] (n_euk) -- node[brlen, above] {0.18} (n_as);
\draw[edge] (n_bac) -- node[brlen, above] {0.22} (n_eb);
\draw[edge] (n_root) -- node[brlen, above, sloped] {0.30} (n_euk);
\draw[edge] (n_root) -- node[brlen, below, sloped] {0.32} (n_bac);

% Vertical connectors (cladogram style)
\draw[edge] (6.5, 4.0) -- (n_as) -- (6.5, 3.2);
\draw[edge] (6.0, 2.0) -- (n_eb) -- (6.0, 1.2);
\draw[edge] (4.5, 3.6) -- (n_euk) -- (4.5, 2.6);
\draw[edge] (4.5, 1.6) -- (n_bac) -- (4.5, 2.6);
\draw[edge] (n_root) -- (2.0, 0.0);

% Bootstrap support labels
\node[bsval] at (6.65, 3.7) {98};
\node[bsval] at (6.15, 1.7) {96};
\node[bsval] at (4.65, 3.0) {87};
\node[bsval] at (4.65, 2.0) {72};

% Internal node annotations
\node[ilabel, anchor=south] at (n_euk) {Eukaryota};
\node[ilabel, anchor=north] at (n_bac) {Bacteria};
\node[ilabel, anchor=east] at (n_root) {LUCA};

% Scale bar
\draw[thick, draw=engblack] (2.0, -1.2) -- (3.0, -1.2);
\node[font=\footnotesize, anchor=north] at (2.5, -1.25)
  {0.1 substitutions per site};

\end{tikzpicture}
\caption[Phylogenetic tree]{Maximum-likelihood phylogeny inferred
from a 1.5~kb 16S/18S rRNA alignment across five taxa, with branch
lengths in substitutions per site. Bootstrap support values (red,
on internal edges) report the fraction of 1000 nonparametric
bootstrap replicates that recover the same partition. Both
bacterial branches and both eukaryote branches receive strong
support; the eukaryote-bacterial split receives moderate support
(72) at this resolution. A branch length is not a calendar age; it
is a substitution count along the lineage, and translating it to
time requires a calibrated molecular clock.}
\label{fig:vol06:ch11:phylo-tree}
\end{figure}


\subsection{Maximum-likelihood and Bayesian methods}

Maximum-likelihood methods evaluate the likelihood
$P(\text{alignment} \mid \text{tree}, \text{substitution model})$ using
Felsenstein's pruning algorithm \cite{paper:v6c11-felsenstein-phylogeny-1981},
which marginalises the unknown ancestral states recursively from the
leaves toward the root in $O(n L)$ time per likelihood evaluation,
where $n$ is the number of taxa and $L$ is the alignment length. The
tree topology is searched by stochastic heuristics (nearest-neighbour
interchange, subtree pruning and regrafting) because the topology
space has $(2n-5)!!$ unrooted topologies, exploding super-exponentially.

Bayesian methods (MrBayes, BEAST) sample tree topologies, branch
lengths, and substitution-model parameters from the posterior
distribution by Markov chain Monte Carlo. The posterior summarises the
uncertainty in all three at once. The cost is computational: a
modest-size analysis (50 taxa, 1500 sites) runs for hours to days; a
large analysis runs for weeks on a cluster.

\subsection{Evolutionary engineering}

The phylogenetic apparatus extends from descriptive to prescriptive
in two engineering contexts. Directed evolution applies iterative
mutation, selection, and recombination to a starting protein or
nucleic acid, navigating the fitness landscape toward a target
function; the phylogenetic vocabulary supplies the population-genetic
intuition for how the population moves under selection. Ancestral
sequence reconstruction infers the most-likely sequence of an
ancestral protein from its descendants, then synthesises and
characterises the ancestor in the laboratory; reconstructed ancestors
often have higher thermostability than their modern descendants and
have been used to seed industrial enzyme-engineering campaigns.

\section{ML for biology: graph networks, foundation models}\pathtag{standard}

\subsection{The supervised learning template in biology}

Most biological-ML tasks fit a supervised template: features extracted
from a biological measurement, a target variable (cell type, disease
state, binding affinity, structural property), and a regression or
classification model. The differences from non-biological ML are not
algorithmic; they are statistical. Biological data are observational,
high-dimensional, low in independent replicates, structured by
phylogeny or cohort, and contaminated by batch effects. The same
neural network that performs well on ImageNet may overfit a 100-sample
omics cohort in an afternoon.

\subsection{Graph neural networks for molecular biology}

Molecules and biological networks are naturally represented as graphs:
nodes are atoms, residues, genes, or proteins; edges are bonds,
interactions, or correlations. Graph neural networks aggregate
information from a node's neighbourhood through learned message-passing
operations, producing per-node or per-graph embeddings suitable for
downstream prediction. Applications include molecular-property
prediction (toxicity, solubility, binding affinity
\cite{paper:v6c11-stark-graphnn-2022}), protein-protein interaction
prediction, and drug-target prediction.

The message-passing operation at layer $\ell$ updates the node
embedding $h_v^{(\ell)}$ as
\[
h_v^{(\ell+1)} = \phi\!\left( h_v^{(\ell)}, \bigoplus_{u \in \mathcal{N}(v)}
\psi(h_u^{(\ell)}, h_v^{(\ell)}, e_{uv}) \right),
\]
where $\mathcal{N}(v)$ is the neighbourhood of $v$, $e_{uv}$ is the edge
feature, $\psi$ is a message function, $\bigoplus$ is a permutation-
invariant aggregator (sum, mean, max), and $\phi$ is an update function.
The number of layers controls the receptive field; for molecular
graphs, three to six layers typically suffice because the relevant
physical interactions are local.

\subsection{Foundation models for sequences}

Protein language models (ESM \cite{paper:v6c11-lin-esm2-2023}, ProtTrans)
and nucleic-acid language models (DNABERT, Nucleotide Transformer)
follow the foundation-model template of large-scale pre-training on
unlabelled sequence data followed by task-specific fine-tuning. ESM-2
is trained on $\sim 10^{8}$ UniRef protein sequences by masked-residue
prediction. The learned embeddings carry residue-level structural
information that transfers to downstream tasks (secondary-structure
prediction, contact prediction, function annotation, variant-effect
prediction) without explicit retraining of the backbone.

The MSA Transformer \cite{paper:v6c11-rao-msa-transformer-2021}
extends the architecture to attend over the rows and columns of a
multiple-sequence alignment, capturing co-evolution explicitly. Its
output is the architectural ancestor of AlphaFold's Evoformer; both
are attention-over-alignments models trained on different objectives.

\subsection{Practical considerations}

Three operational issues are specific to biological ML. Hyperparameter
tuning must use proper cross-validation that respects the cohort
structure (patient-level, site-level, or cell-line-level folds, not
sample-level folds). Reporting must include external-validation
performance, that is, performance on a cohort independent of the
training distribution; internal cross-validation is a development
metric and not a release metric. Confidence calibration matters more
than peak accuracy: an omics classifier used to triage patients must
report a well-calibrated probability, because an over-confident
prediction at the wrong threshold becomes the failure case described
in the next section.

\section{Failure: data leakage in biological ML, batch effects, overfit-to-cohort}\pathtag{core}

\subsection{Three failure mechanisms with a common signature}

Biological ML pipelines fail in three related ways. Data leakage admits
information about the held-out set into the training process; the
model learns to recognise the leakage and reports the recognition as
generalisation. Batch effects produce systematic differences between
samples processed in different runs, machines, sites, or days; if
batch is correlated with the outcome of interest, the model learns
batch instead of biology. Cohort overfit fits a model that performs
well on the training cohort and fails on any other cohort because the
training cohort's idiosyncrasies (genetic background, demographics,
treatment history, ascertainment bias) are not shared by other
cohorts.

All three produce the same surface symptom: optimistic internal
performance, failure on external validation, sometimes accompanied by
a paper, a press release, or a clinical trial. The internal symptom
is harder to read because it is sample-by-sample correct; the failure
sits in the cohort, not in any individual case.

\subsection{Data leakage}

\Cref{fig:vol06:ch11:ml-leakage} sketches the four canonical pathways
of leakage. Patient-level versus sample-level splitting: if two
samples from one patient end up in the training and test sets
respectively, the model learns the patient (the genotype, the imaging
artefact, the institutional preference) rather than the disease, and
reports the same-patient match as out-of-sample performance.
Preprocessing on the full cohort: normalisation parameters, batch-
correction vectors, feature-selection rankings, and PCA loadings that
use information from the held-out samples leak that information into
the training pipeline. Hyperparameter tuning on the test set: model
selection or threshold tuning on the test set converts the test set
into a second training set, and the final reported performance is
optimistic by the amount of test-set adaptation. Same-source
contamination: training and test sharing imaging sites, sequencing
centres, or scanner instruments allows the model to learn the
site-specific artefact instead of the biological signal.

% Figure: ML-leakage schematic showing the four canonical pathways
% by which test data contaminates training in biological ML.
\begin{figure}[htbp]
\centering
\begin{tikzpicture}[
  pool/.style={draw=engblack, rounded corners=2pt, minimum width=2.8cm,
    minimum height=1.0cm, align=center, font=\small},
  leak/.style={->, thick, draw=engred!80, dashed},
  ok/.style={->, thick, draw=engblack!70},
  bad/.style={fill=engred!10},
  good/.style={fill=englime!10},
]

% Top row: the intended pipeline
\node[pool, good] (raw) at (0, 3) {Raw cohort};
\node[pool, good] (split) at (3.5, 3) {Patient-level\\hold-out split};
\node[pool, good] (train) at (7.5, 3.8) {Training set\\(80\%)};
\node[pool, good] (test) at (7.5, 2.2) {Test set\\(20\%)};
\node[pool, good] (mod) at (11.5, 3.8) {Train model};
\node[pool, good] (eval) at (11.5, 2.2) {Evaluate};

\draw[ok] (raw) -- (split);
\draw[ok] (split) -- (train);
\draw[ok] (split) -- (test);
\draw[ok] (train) -- (mod);
\draw[ok] (test) -- (eval);
\draw[ok] (mod.south) -| (eval.north);

% Below: four leakage paths
\node[font=\small\bfseries, anchor=west] at (0, 0.8)
  {Four leakage pathways:};

\node[pool, bad] (l1) at (1.7, -0.7) {(1) Sample-level\\not patient-level\\split};
\node[pool, bad] (l2) at (5.2, -0.7) {(2) Preprocessing\\on full cohort\\(normalisation,\\feature selection)};
\node[pool, bad] (l3) at (8.7, -0.7) {(3) Hyperparameter\\tuning on\\test set};
\node[pool, bad] (l4) at (12.2, -0.7) {(4) Same-source\\contamination\\(images, sites,\\scanners)};

% Show paths leaking into the test set
\draw[leak] (l1.north) to[bend right=15] (test.south);
\draw[leak] (l2.north) to[bend left=10] (test.south);
\draw[leak] (l3.north) to[bend left=20] (test.south);
\draw[leak] (l4.north) to[bend left=30] (test.south);

\end{tikzpicture}
\caption[Data-leakage pathways in biological ML]{Four pathways by
which test information leaks into a biological-ML pipeline. (1)
Splitting the cohort at the sample level when multiple samples come
from the same patient (the same tumour resected at two timepoints,
two cells from the same individual): the model learns
patient-specific signatures and reports them as generalisation. (2)
Preprocessing (normalisation, batch correction, feature selection)
fitted on the full cohort uses information from the held-out
samples to set parameters that the model then exploits. (3)
Hyperparameter selection by performance on what the engineer calls
the test set converts the test set into a second training set. (4)
Same-source contamination: when training and test share imaging
sites, sequencing centres, or scanner instruments, the model can
learn the site-specific artefact instead of the biological signal.
The Duke Potti case combined pathways 1, 2, and 3 with reporting
errors; the COVID-19 imaging review of 415 papers found pathway 4
endemic.}
\label{fig:vol06:ch11:ml-leakage}
\end{figure}


The Duke Potti case \cite{paper:v6c11-baggerly-coombes-2009,
paper:v6c11-iom-omics-2012} combined pathways one, two, and three with
additional reporting errors. The Potti group at Duke University Medical
Center published genomic predictors of chemotherapy response between
2006 and 2007, used those predictors to assign treatment in three
clinical trials between 2007 and 2010, and saw the predictors
implemented in commercial diagnostic products. Baggerly and Coombes at
MD Anderson reanalysed the published predictors beginning in 2007,
finding swapped sensitive-and-resistant labels on the training data,
off-by-one sample index shifts, overlapping training and test sets,
and an inability to regenerate the published predictors from the
supplied supplementary materials. The clinical trials were suspended
in 2010, and several of the underlying papers were retracted between
2010 and 2015. The Institute of Medicine's 2012
report \emph{Evolution of Translational Omics}, commissioned by NIH
after the case, codified a recommended framework for omics-based
predictor evaluation: locked computational predictors, independent
held-out validation cohorts, predefined performance thresholds, and
institutional-review-board approval that explicitly addresses the
predictor's validation status \cite[][IOM Recommendations
1--4]{paper:v6c11-iom-omics-2012}.

\subsection{Batch effects}

Batch effects are systematic differences between groups of samples
that depend on processing factors rather than on the biological factor
of interest. Sources include sequencing run, library-preparation
reagent lot, instrument, technician, time of day, freezer storage time,
and laboratory of origin. The Leek et al. 2010 review
\cite{paper:v6c11-leek-batch-2010} surveys the canonical mechanisms
across genomics.

\Cref{fig:vol06:ch11:batch-effect} shows a paradigmatic case. In the
left panel, PC1 separates samples by laboratory of origin and the
disease signal (healthy versus tumour) sits on PC2. A classifier
trained on PC1 will appear to perform well on within-laboratory tests
because PC1 captures real variance, but will misclassify when applied
to a new laboratory because the new laboratory's PC1 has a different
sign. The right panel shows the result after ComBat correction
\cite{paper:v6c11-johnson-combat-2007}: PC1 now captures the disease
signal, and healthy and tumour clusters align across laboratories.

% Figure: Batch-effect contamination demonstration in PCA space
% before and after ComBat correction.
\begin{figure}[htbp]
\centering
\begin{tikzpicture}
% Left panel: before correction
\begin{scope}
\begin{axis}[
  width=7cm, height=6cm,
  name=before,
  title={\textbf{Before correction}: PC1 driven by batch},
  title style={font=\small},
  xlabel={PC1 (52\% var.)},
  ylabel={PC2 (14\% var.)},
  xmin=-5, xmax=5,
  ymin=-3, ymax=3,
  grid=both,
  major grid style={draw=engblack!15},
  ticks=none,
  legend pos=south east,
  legend style={font=\footnotesize, draw=engblack!40},
]

% Batch 1 (lab A): healthy
\addplot[only marks, mark=*, mark size=2pt, color=engblue]
  table {
  x y
  -3.2 1.5
  -3.0 1.0
  -3.5 1.8
  -2.8 0.5
  -3.4 1.2
  -3.1 1.6
  -2.9 0.8
  -3.3 1.4
  };
\addlegendentry{Lab A, healthy}

% Batch 1 (lab A): tumour
\addplot[only marks, mark=square*, mark size=2pt, color=engblue]
  table {
  x y
  -2.5 -1.0
  -2.7 -0.6
  -2.3 -1.4
  -2.6 -0.8
  -2.8 -1.2
  -2.4 -0.9
  };
\addlegendentry{Lab A, tumour}

% Batch 2 (lab B): healthy
\addplot[only marks, mark=*, mark size=2pt, color=engamber]
  table {
  x y
  3.0 1.0
  3.2 1.4
  2.8 0.6
  3.4 1.2
  3.1 1.5
  2.9 0.9
  3.3 1.1
  };
\addlegendentry{Lab B, healthy}

% Batch 2 (lab B): tumour
\addplot[only marks, mark=square*, mark size=2pt, color=engamber]
  table {
  x y
  3.5 -1.2
  3.7 -0.8
  3.4 -1.5
  3.6 -1.0
  3.8 -1.3
  3.5 -0.7
  };
\addlegendentry{Lab B, tumour}

\end{axis}
\end{scope}

% Right panel: after ComBat
\begin{scope}[xshift=8cm]
\begin{axis}[
  width=7cm, height=6cm,
  title={\textbf{After ComBat}: PC1 driven by disease},
  title style={font=\small},
  xlabel={PC1 (38\% var.)},
  ylabel={PC2 (10\% var.)},
  xmin=-4, xmax=4,
  ymin=-3, ymax=3,
  grid=both,
  major grid style={draw=engblack!15},
  ticks=none,
  legend pos=south east,
  legend style={font=\footnotesize, draw=engblack!40},
]

% Healthy from both labs cluster together
\addplot[only marks, mark=*, mark size=2pt, color=engblue]
  table {
  x y
  -2.4 0.5
  -2.6 0.8
  -2.2 1.0
  -2.5 0.7
  -2.7 0.6
  -2.3 0.9
  -2.5 0.4
  -2.6 1.1
  };

\addplot[only marks, mark=*, mark size=2pt, color=engamber]
  table {
  x y
  -2.0 0.3
  -2.2 0.7
  -1.8 1.0
  -2.1 0.5
  -2.3 0.8
  -1.9 0.6
  -2.1 0.9
  };

% Tumour from both labs cluster together
\addplot[only marks, mark=square*, mark size=2pt, color=engblue]
  table {
  x y
  2.5 -0.6
  2.7 -1.0
  2.3 -0.4
  2.6 -0.8
  2.8 -1.1
  2.4 -0.5
  };

\addplot[only marks, mark=square*, mark size=2pt, color=engamber]
  table {
  x y
  2.2 -0.7
  2.4 -1.1
  2.0 -0.5
  2.3 -0.9
  2.5 -1.2
  2.1 -0.6
  };

\end{axis}
\end{scope}

\end{tikzpicture}
\caption[Batch effect before and after correction]{Principal-component
projection of a simulated dual-batch expression dataset, with
healthy (circle) and tumour (square) samples drawn from two labs
(blue and amber). Left: before correction, PC1 separates lab A from
lab B; the disease signal sits on PC2 and is dominated by
batch-of-origin variance. A classifier trained on PC1 will appear
to perform well at predicting disease within each lab but will fail
when applied to a new lab. Right: after empirical-Bayes batch
correction (ComBat), the batch axis is suppressed and the disease
signal appears on PC1; healthy and tumour clusters now align across
labs. ComBat removes mean and variance shifts associated with batch
under the model assumption that batch and disease are not
collinear. If the experimental design is collinear (e.g., all lab A
samples are healthy and all lab B samples are tumour), ComBat
removes the disease signal along with the batch signal.}
\label{fig:vol06:ch11:batch-effect}
\end{figure}


ComBat is an empirical-Bayes adjustment that fits per-feature, per-
batch location and scale parameters and shrinks them toward a pooled
prior. The model is
\[
y_{ijg} = \alpha_{g} + X\beta_{g} + \gamma_{ig} + \delta_{ig}\,\epsilon_{ijg},
\]
where $y_{ijg}$ is the expression of feature $g$ in sample $j$ of
batch $i$, $\alpha_{g}$ is the overall mean, $X\beta_{g}$ is the
biological-condition effect, $\gamma_{ig}$ is the additive batch
effect, $\delta_{ig}$ is the multiplicative batch scale, and
$\epsilon_{ijg}$ is the residual. Empirical-Bayes estimation borrows
strength across features within a batch. The adjusted measurement
$y_{ijg}^{\ast}$ removes the estimated $\gamma_{ig}$ and rescales by
the estimated $\delta_{ig}$.

The critical assumption is that batch and condition are not collinear.
If laboratory A processed only healthy samples and laboratory B
processed only tumour samples, ComBat cannot tell whether PC1 is
laboratory or disease and will remove the disease signal as it removes
the laboratory signal. The experimental-design rule that follows is
that every batch must contain samples from every condition, in
proportions that preserve identifiability. The rule is honoured more
in the breach than in the observance.

\subsection{Overfit-to-cohort}

A model overfits its cohort when its training distribution differs
systematically from the deployment distribution and the differences
are not captured by the features the model sees. Wong et al.
\cite{paper:v6c11-wong-sepsis-2021} externally validated the Epic
Sepsis Model, a widely deployed proprietary sepsis-prediction model,
at the University of Michigan health system and found that the model's
area under the ROC was 0.63 (well below the 0.76 to 0.83 claimed by the
vendor) and that the model missed 67$\%$ of true sepsis cases at the
deployed threshold. The mechanism was a training-distribution mismatch:
the model had been trained on a different patient mix with different
diagnostic-coding practices, and the deployment institution's coding
practice was systematically different from the training institution's.
The model did not see the relevant difference because the relevant
difference was not in any of the features it used.

The COVID-19 imaging-AI review of Roberts et al.
\cite{paper:v6c11-roberts-covid-ml-2021} systematically evaluated $415$
machine-learning papers that claimed to detect or prognosticate
COVID-19 from chest radiographs or CT scans, published between January
and October 2020. None of the $415$ models were judged clinically
usable. The most common failure mode was same-source contamination:
COVID-positive and COVID-negative images came from different hospitals,
and the models learned the hospital instead of the disease. Other
failure modes included sample-level rather than patient-level splits,
heavy class imbalance with no calibrated handling, and lack of
external-cohort validation. The paper is the largest single audit of
biological-ML failure modes in the post-2020 literature and is the
canonical reference for the operational rules that govern
biological-ML reporting.

The Google Flu Trends episode \cite{paper:v6c11-lazer-google-flu-2014}
illustrates the cohort-overfit mechanism on a population scale: a
predictor trained on the correlation between flu-related Google search
queries and CDC flu reports during 2003 to 2008 performed well in
hold-out periods through 2010 and then over-predicted flu prevalence
by a factor of two during the 2012-2013 season because the search-
behaviour distribution had shifted. The relevant difference was not in
the model's features.

\begin{principle}[Biological structure is the signal; the engineer's discipline preserves it across the pipeline]
A biological dataset arrives with structure: gene families, pathways,
patient cohorts, sequencing batches, instrument lineages,
phylogenetic relationships, and tissue contexts. The structure is
the signal that the algorithm is meant to extract; the same structure
is the route by which an undisciplined pipeline confuses cohort for
biology, batch for condition, site for disease, and patient for
population. The engineer's task is to design the pipeline so that
each step preserves the biological signal and respects the
non-biological structure as a constraint. The Duke Potti, Epic Sepsis,
and COVID-19 imaging failures all violated this rule at a
specific identifiable step: training and test from the same patient,
preprocessing fitted on the full cohort, hyperparameter tuning on the
test set, and same-source contamination respectively.
\end{principle}

\subsection{The operational rules}

The post-Duke-Potti operational rules for biological ML are explicit.
Lock the computational predictor before validation: every
preprocessing parameter, every feature selection, every hyperparameter,
and every threshold is fixed before the validation cohort is touched.
Validate on a cohort generated independently of the discovery cohort,
ideally at a different institution and on a different patient
population. Pre-register the analysis pipeline and the performance
thresholds. Report the cohort metadata fully, including the inclusion
and exclusion criteria, the sample-collection procedures, the
processing batches, and the per-batch sample composition. Use
patient-level (or site-level, or cell-line-level) cross-validation
folds, not sample-level folds. Quote calibration as a first-class
performance metric, not only discrimination. Treat the FDA software-
as-a-medical-device pathway as the regulatory binding for predictors
intended for clinical use. The Institute of Medicine recommendations
1 through 4 \cite{paper:v6c11-iom-omics-2012} codify these rules for
the United States NIH-funded translational-omics community; the FDA
NGS guidance \cite{std:v6c11-fda-ngs-2018} codifies them for
diagnostic NGS panels.

\section{Project}\pathtag{core}

\begin{project}[Align two bacterial genomes and annotate variants]
\textbf{Goal.}
Take two short bacterial-genome fragments, perform global pairwise
alignment, identify the variant positions, classify each variant by
the kind of change it represents, and produce a one-page annotation
summary suitable for a biology collaborator.

\textbf{Phases.}
The project runs in five phases.

Phase 1: Alignment. Use the supplied pair
\repopath{volumes/06-life/ch11-bioinformatics/data/genome_a.fasta} and
\repopath{volumes/06-life/ch11-bioinformatics/data/genome_b.fasta}
($\approx 5$~kb each, synthetic). Run
\repopath{volumes/06-life/ch11-bioinformatics/code/needleman_wunsch.py}
on the two sequences with match $+1$, mismatch $-1$, gap $-1$. Save the
output alignment to a file. Report the alignment score and the number
of matched, mismatched, and gapped positions.

Phase 2: Variant identification. Walk the alignment column by column.
For each column, classify the type: match, transition (purine-to-purine
or pyrimidine-to-pyrimidine substitution), transversion (purine-to-
pyrimidine substitution or vice versa), insertion (gap in genome A),
deletion (gap in genome B). Produce a list of variants, each with its
position in the alignment, its REF and ALT alleles, and its type.

Phase 3: Filter and quality. Construct a hypothetical VCF for the
variants you identified, with synthetic QUAL, DP, FS, and MQ values
drawn from a plausible distribution. Run
\repopath{volumes/06-life/ch11-bioinformatics/code/variant_filter.py}
on it. Report which variants pass and which fail each filter, and
explain how the synthetic distribution influenced the pass rate.

Phase 4: Annotation. For each PASS variant, look up its position in
the (synthetic) gene-model GFF that you write by hand for the
fragment. Annotate each variant as intergenic, synonymous, missense
(specify amino-acid change), nonsense, or frameshift. Provide a
codon-table reference.

Phase 5: Summary. Write a one-page report that contains: the alignment
score and identity, the variant count by type, the filter pass rate,
the functional annotations, and a one-paragraph honest statement of
what this analysis can and cannot support. The honest statement is
the load-bearing deliverable.

\textbf{Deliverables.}
A working alignment file, a tab-delimited variant table, a filtered
VCF, a hand-written gene-model GFF, an annotation table, and a
one-page summary report.

\textbf{Success criteria.}
The alignment is correct (the score and identity match a reference
calculation). The variant table contains every column-level
disagreement between the two genomes (no missing variants and no
spurious matches reported as variants). The filter pipeline rejects
the variants whose QUAL, DP, FS, or MQ fall outside the thresholds.
The annotation is internally consistent with the gene-model GFF. The
one-page summary states clearly what the analysis supports and what
it does not (it does not, for instance, support a statement about
biological consequence, because the gene model is synthetic).

\textbf{Common failure modes.}
Reporting alignment identity as the fraction of identical positions
in the aligned region without correcting for the gaps. Reporting
variants only at mismatch positions and forgetting indels. Setting
filter thresholds by what makes the most variants pass rather than by
what is defensible from the per-call statistics. Annotating
synonymous-versus-missense without first identifying the reading
frame. Writing a summary that claims more than the analysis supports.

\textbf{Worked example pointer.}
The synthetic pair $A$ and $B$ in
\repopath{volumes/06-life/ch11-bioinformatics/data/pair_example.fasta}
is a 50-nucleotide pair with about 4 substitutions and one short
indel. The Needleman-Wunsch traceback recovers the unique optimal
alignment for this pair; the variant table contains a small number of
SNVs (mix of transitions and transversions) and one indel. Use this
pair to verify your pipeline before running on the 5~kb pair.
\end{project}

\section{Exercises}

\subsection*{Calculation}\pathtag{standard}

\begin{exercise}
Compute the Needleman-Wunsch optimal global alignment of the
sequences AAGCT and AGCT using match $+1$, mismatch $-1$, gap $-1$.
Write out the full $6 \times 5$ DP matrix and one optimal alignment.

\paragraph{Solution sketch.}
Initialise: $M[0,0]=0$, first row $0,-1,-2,-3,-4$, first column
$0,-1,-2,-3,-4,-5$. Fill row by row using the recurrence:
$M[1,1]=\max(0+1, -1-1, -1-1)=1$ (match A-A).
$M[1,2]=\max(-1-1, 1-1, -2-1)=0$ etc.
The final matrix has $M[5,4]=3$. One optimal alignment is
AAGCT vs A-GCT with a single gap aligning the second A in the
first sequence to a gap; the alignment has $4$ matches and one gap
($4 - 1 = 3$).
\end{exercise}

\begin{exercise}
A BLAST search of a $300$-residue query against a $10^{9}$-residue
protein database returns a top hit with raw score $S = 130$ bits
under BLOSUM62 with $K = 0.1$ and $\lambda = 0.3$ per bit. Compute the
E-value and the normalised bit score. Is the hit significant at the
$E = 0.01$ threshold?

\paragraph{Solution sketch.}
$E = K m n e^{-\lambda S} = 0.1 \times 300 \times 10^{9} \times
e^{-0.3 \times 130} = 3 \times 10^{10} \times e^{-39} \approx 3 \times
10^{10} \times 1.2 \times 10^{-17} \approx 3.6 \times 10^{-7}$.
Normalised bit score $S' = (\lambda S - \ln K)/\ln 2 = (39 + 2.30)/0.693
\approx 59.6$ bits. With $E \approx 4 \times 10^{-7} \ll 0.01$, the
hit is significant. Confirmation: $m n / 2^{S'} = 3 \times 10^{11} /
2^{59.6} \approx 3 \times 10^{11} / 8.7 \times 10^{17} \approx 3.5
\times 10^{-7}$, agreeing.
\end{exercise}

\begin{exercise}
The Phred quality string for the first 10 bases of a read reads
\texttt{!"\#\$\%\&'()*} (Sanger encoding, ASCII offset 33). Compute the
Phred scores and the per-base error probabilities for each of the 10
positions, and decide whether the read passes a $Q_{\mathrm{min}}=20$
filter at all 10 positions.

\paragraph{Solution sketch.}
Subtract 33 from each ASCII code: ! is $33$ so $Q=0$ ($p=1$); " is
$34$ so $Q=1$ ($p\approx 0.79$); \# is $35$ so $Q=2$ ($p\approx 0.63$);
\$ is $36$ so $Q=3$; through $*$ is $42$ so $Q=9$. Per-base error
probabilities $p = 10^{-Q/10}$. None of the ten positions reaches
$Q=20$ (the minimum $p$ is $10^{-0.9} \approx 0.126$ at position 10),
so the entire prefix fails the $Q_{\mathrm{min}}=20$ filter; a trimmer
would discard the read or trim the prefix to its first position with
$Q\geq 20$.
\end{exercise}

\begin{exercise}
A synthetic gene of length $L=900$~bp has GC content $0.55$. Compute
the expected number of \texttt{GAATTC} sites (the EcoRI recognition
sequence) in a uniform-random sequence of the same length and GC
content. Assume positions are independent.

\paragraph{Solution sketch.}
For GC content $g=0.55$, $\Pr(\text{G}) = \Pr(\text{C}) = g/2 = 0.275$
and $\Pr(\text{A}) = \Pr(\text{T}) = (1-g)/2 = 0.225$. The motif
\texttt{GAATTC} has probability
$p = 0.275 \times 0.225 \times 0.225 \times 0.225 \times 0.225 \times
0.275 = 0.275^{2} \times 0.225^{4} \approx 0.0756 \times 0.00256
\approx 1.94 \times 10^{-4}$. Expected sites in a sequence of length
$L=900$ is $(L - 6 + 1) \times p \approx 895 \times 1.94 \times
10^{-4} \approx 0.17$. The reasonable interpretation: under the null,
this short region is expected to contain less than one EcoRI site;
observing two or more would be moderately surprising.
\end{exercise}

\begin{exercise}
Given the $5 \times 5$ distance matrix in
\repopath{volumes/06-life/ch11-bioinformatics/code/phylo_distance.py}
(species Aspergillus, Saccharomyces, E.~coli, B.~subtilis,
Halobacterium with pairwise distances 0.13 to 0.74), execute one step
of UPGMA by hand: identify the closest pair, compute the height of
the join, write the partial Newick string after the merge, and write
the updated $4 \times 4$ distance matrix.

\paragraph{Solution sketch.}
Closest pair: (Aspergillus, Saccharomyces) at distance $0.13$. Height
of join: $0.13/2 = 0.065$. Each branch length from the leaf to the new
node is $0.065 - 0 = 0.065$ (both leaves are at height 0). Partial
Newick: (Aspergillus:0.065,Saccharomyces:0.065). Updated distance
from the new cluster $(A,S)$ to the others: each is the mean of the
two leaf distances. To E.~coli: $(0.62+0.59)/2 = 0.605$.
To B.~subtilis: $(0.58+0.61)/2 = 0.595$. To Halobacterium:
$(0.72+0.74)/2 = 0.73$. The other off-diagonal distances are
unchanged. The next UPGMA step joins E.~coli and B.~subtilis at
distance $0.23$ (height $0.115$).
\end{exercise}

\subsection*{Derivation}\pathtag{standard}

\begin{exercise}
Derive the Karlin-Altschul E-value formula
$E = K m n e^{-\lambda S}$ from the extreme-value distribution for
maxima of i.i.d. random variables. State the assumptions that go into
$K$ and $\lambda$.

\paragraph{Solution sketch.}
Consider all ungapped local alignments between a query of length $m$
and a database of length $n$. Under the null model (independent random
residues drawn from a fixed background), the number of distinct local
alignments is proportional to $m n$. Each alignment's score is the sum
of $L$ independent substitution scores, where $L$ is the alignment
length. For score systems with negative expected score per position
(required to make the maximum non-degenerate), the local-alignment
maximum follows the Gumbel extreme-value distribution: $\Pr(S_{\max} >
x) \approx 1 - \exp(-K m n e^{-\lambda x})$ for large $m, n$. The
expected number of distinct hits with score above $S$ is the
right-tail of this distribution, $E \approx K m n e^{-\lambda S}$.
The parameter $\lambda$ is the unique positive root of $\sum_{a,b}
p_{a} p_{b} e^{\lambda s(a,b)} = 1$, where $p_{a}$ is the residue
frequency. $K$ is a correction factor that depends on the score
distribution. The derivation requires: ungapped alignments, fixed
background composition, score system with negative expected score per
position, and large $m, n$.
\end{exercise}

\begin{exercise}
Derive the Jukes-Cantor distance-correction formula
$d = -(3/4) \ln(1 - (4/3)p)$ from a continuous-time Markov chain on
the four nucleotides with uniform substitution rate $\mu$.

\paragraph{Solution sketch.}
Under Jukes-Cantor, the rate matrix has $\mu$ in each off-diagonal
position and $-3\mu$ on the diagonal. The transition probability over
time $t$ has $\Pr(\text{same}\mid t) = 1/4 + (3/4) e^{-4\mu t}$ and
$\Pr(\text{different} \mid t) = 3/4 - (3/4) e^{-4\mu t}$. The observed
mismatch fraction $p = 3/4 - (3/4) e^{-4\mu t}$. Solve for $4\mu t$:
$e^{-4\mu t} = 1 - (4/3) p$, so $4\mu t = -\ln(1 - (4/3) p)$. The
substitution distance is the expected number of substitutions per
site, $d = 3\mu t = -(3/4) \ln(1 - (4/3) p)$. The formula breaks down
when $p \to 3/4$ (saturation): the logarithm diverges, reflecting that
beyond this point the observed mismatch fraction is consistent with
arbitrarily large numbers of substitutions.
\end{exercise}

\begin{exercise}
Derive the chemostat washout condition for a Monod growth model in
continuous culture. (This exercise draws on the bioreactor balance
from Chapter 10; here, the derivation is the engineering parallel of
the BLAST-search statistic: a threshold above which the system
behaviour changes qualitatively.)

\paragraph{Solution sketch.}
Let $X$ be biomass concentration, $S$ substrate concentration, $D$
dilution rate, $S_{f}$ feed substrate concentration. The chemostat
balance is $dX/dt = (\mu(S) - D) X$, $dS/dt = D(S_{f} - S) -
\mu(S) X / Y$, with Monod kinetics $\mu(S) = \mu_{\max} S/(K_{S} +
S)$. At nonzero steady state, $\mu(S^{\ast}) = D$, so $S^{\ast} =
K_{S} D / (\mu_{\max} - D)$. The steady-state $S^{\ast}$ is well
defined and positive only if $D < \mu_{\max}$. As $D \to \mu_{\max}$,
$S^{\ast} \to \infty$, which is impossible because $S^{\ast} \leq
S_{f}$; the practical washout limit is $D_{\max} = \mu_{\max} S_{f}
/ (K_{S} + S_{f})$. For $D > D_{\max}$, the only steady state is
$X = 0, S = S_{f}$ (sterile pipe). The bioinformatics parallel: the
washout condition is the equivalent of the BLAST E-value crossing the
reporting threshold; both are single inequalities that switch the
system from "useful" to "not useful" qualitatively.
\end{exercise}

\begin{exercise}
Show that maximum-likelihood phylogeny inference under a time-reversible
substitution model is independent of the choice of root for an unrooted
binary tree. (This is the "pulley principle" of Felsenstein.)

\paragraph{Solution sketch.}
Under a time-reversible model, $\pi_{a} P_{ab}(t) = \pi_{b} P_{ba}(t)$
for stationary frequencies $\pi$. The likelihood at the root is
$\sum_{r} \pi_{r} \prod_{c} L_{c}(r)$ where $L_{c}(r)$ is the
conditional likelihood of subtree $c$ given root state $r$. Moving the
root from one position on an edge to another redistributes the edge
length between the two sides of the root; by time-reversibility, the
likelihood is unchanged because the product $\pi_{r} P_{rs}(t)$ on the
new root-adjacent edges equals the same product computed at the old
root. The pulley principle holds across edges of the tree as well: the
root may be placed anywhere on the unrooted tree without changing the
likelihood. The consequence for software design is that tree-likelihood
computation may use any edge as the root, which Felsenstein's pruning
algorithm exploits for efficiency.
\end{exercise}

\subsection*{Estimation}\pathtag{standard}

\begin{exercise}
Estimate the storage cost for one year of raw FASTQ output from a
single Illumina NovaSeq 6000 instrument operated at full capacity.
Use commercial cloud-storage pricing (current as of 2024).

\paragraph{Solution sketch.}
NovaSeq 6000 with S4 flow cell produces approximately $3$~TB
of compressed FASTQ per run (about $20$~billion reads at $150$~bp
paired end). One full run takes $\approx 44$~h; with maintenance and
turnover, a single instrument can complete $\approx 160$ runs per year,
producing $\approx 480$~TB. At AWS S3 standard pricing of approximately
\$0.023 per GB-month (current as of 2024), one year of storage costs
$480 \times 10^{3} \times 0.023 \times 12 = $~\$$132{,}000$ for one
year of one instrument's output, before egress or compute. Compression
to CRAM with the reference genome reduces the storage by a factor of
$\approx 3$, cutting the bill to $\approx$~\$$45{,}000$. Many large
sequencing centres operate $20$ or more such instruments, putting the
annual storage cost into the millions of dollars per institution;
this is the operational driver behind the move to archive only the
called variants for routine cohort projects.
\end{exercise}

\begin{exercise}
Estimate the number of UniProt entries (current as of 2024) and the
fraction that are computationally annotated versus manually curated.

\paragraph{Solution sketch.}
UniProtKB has two sections: Swiss-Prot (manually curated) and TrEMBL
(computationally annotated). Current as of 2024, Swiss-Prot has
approximately $570{,}000$ entries; TrEMBL has approximately $250 \times
10^{6}$ entries. The manual fraction is therefore approximately
$2 \times 10^{-3}$, two parts per thousand. The growth gap is the
operational driver behind structure-prediction databases: a manually
curated PDB structure costs months of effort; AlphaFold and ESM
predictions fill the TrEMBL space at compute cost. The corresponding
curation discipline is to read which database an annotation came from
before quoting it.
\end{exercise}

\begin{exercise}
Estimate the number of clinically actionable variants typically
reported from a single human whole-genome sequence in a research-grade
clinical pipeline.

\paragraph{Solution sketch.}
A human WGS yields $\sim 4 \times 10^{6}$ raw single-nucleotide
variants. After quality filtering, $\sim 3.5 \times 10^{6}$ pass.
Restricting to coding regions: human coding sequence is $\sim 1\%$
of the genome, so $\sim 3.5 \times 10^{4}$ coding variants. Restricting
to missense, nonsense, frameshift, and canonical splice variants:
$\sim 1.5 \times 10^{4}$ predicted-deleterious variants. Restricting
to known disease-gene panels (ClinVar pathogenic or likely pathogenic
in a curated gene list): $\sim 5$ to $50$ variants. Restricting to
variants relevant to the patient's clinical question: typically $0$ to
$5$. The funnel from $4 \times 10^{6}$ raw variants to $\sim 5$
actionable variants is the per-patient version of the same
information-extraction problem that runs throughout the chapter: most
of the engineering work is filtering, not finding.
\end{exercise}

\begin{exercise}
Estimate the number of single-cell RNA-seq experiments deposited in
the Gene Expression Omnibus (GEO) per year. Estimate the implied
aggregate cost.

\paragraph{Solution sketch.}
GEO deposits in 2023 exceeded $\sim 300{,}000$ submissions across all
assay types. Estimates from the Human Cell Atlas project and other
surveys (current as of 2024) place scRNA-seq submissions at
$\sim 20{,}000$ per year, with a median submission containing
$\sim 50{,}000$ cells. At a per-cell cost of $\$0.10$ to $\$0.50$ (a
rough median that includes reagents, sequencing, and pipeline compute),
the per-submission cost is $\sim \$5{,}000$ to $\$25{,}000$. The
aggregate global investment in deposited scRNA-seq data is therefore
on the order of $\$100$ million to $\$500$ million per year, before
the cost of associated computational analysis. The order-of-magnitude
takeaway: scRNA-seq is now a globally distributed measurement
infrastructure on the same financial scale as a small satellite
constellation.
\end{exercise}

\subsection*{Simulation}\pathtag{standard}

\begin{exercise}
Use
\repopath{volumes/06-life/ch11-bioinformatics/code/blast_evalue.py}
to compute E-values for a single hit of bit score $S' = 40$ against
databases of size $10^{6}, 10^{7}, 10^{8}, 10^{9}, 10^{10}$ residues
with query length $m = 300$. Plot $\log_{10} E$ as a function of
$\log_{10} n$ and verify the slope.

\paragraph{Solution sketch.}
Use $E = m n / 2^{S'}$: at $S' = 40$, $E = 300 n / 2^{40} \approx
2.7 \times 10^{-10} n$. For $n = 10^{6}$ to $10^{10}$, $E$ ranges from
$2.7 \times 10^{-4}$ to $2.7$. The slope of $\log_{10} E$ versus
$\log_{10} n$ is exactly $1$, because $E$ is linear in $n$. The same
bit score that is highly significant against a $10^{6}$-residue
database is borderline against a $10^{10}$-residue database; doubling
the database doubles the E-value. The simulation confirms the linear
scaling and gives the student a graphical anchor for the dependence.
\end{exercise}

\begin{exercise}
Run
\repopath{volumes/06-life/ch11-bioinformatics/code/batch_correct.py}
with the default simulation parameters. Report the correlation of PC1
with batch and with disease before and after ComBat. Then modify the
simulation so that batch and disease are collinear (all batch-0 samples
are healthy and all batch-1 samples are tumour) and run again. Explain
the change.

\paragraph{Solution sketch.}
In the default simulation, batch and disease are crossed, and ComBat
restores the disease signal on PC1. Expected output: $\text{corr(PC1,
batch)}$ moves from $\approx +0.95$ before to $\approx 0$ after;
$\text{corr(PC1, disease)}$ moves from $\approx 0$ to $\approx +0.9$.
In the collinear modification, batch and disease produce the same
column in the design matrix; ComBat removes the batch shift, which
also removes the disease shift, and the post-correction PC1 captures
residual noise. The correlation with disease drops near zero. The
exercise illustrates the experimental-design rule: a batch must
contain samples from every condition, in proportions that preserve
identifiability; ComBat does not rescue a collinear design.
\end{exercise}

\begin{exercise}
Simulate sample-level versus patient-level cross-validation on a
synthetic dataset with $200$ samples drawn from $40$ patients ($5$
samples each). Generate the target as a function of patient identity
(strong) plus measurement noise (weak). Compare the apparent
cross-validation accuracy under sample-level folds (random
$80$/$20$ split) and patient-level folds (random split of patient
IDs). Report the gap and explain.

\paragraph{Solution sketch.}
Under sample-level folds, every test sample has training samples from
the same patient; the model learns the patient identifier and reports
high accuracy ($\approx 95\%$ if the patient effect is strong).
Under patient-level folds, no test patient has any training sample;
the model can only generalise via features that share between
patients, and the accuracy drops to chance ($\approx 50\%$ if the
patient effect is the only signal). The gap of $\approx 45$
percentage points is the leakage. This is the canonical sample-level
versus patient-level split error from
\Cref{fig:vol06:ch11:ml-leakage}.
\end{exercise}

\begin{exercise}
Simulate the Karlin-Altschul E-value distribution by generating
$10^{6}$ pairs of random protein sequences of length $100$ each from
the BLOSUM62 background, computing the maximum ungapped alignment
score for each pair, and comparing the resulting empirical
distribution to the Gumbel prediction.

\paragraph{Solution sketch.}
Draw residues independently from the BLOSUM62 background frequencies;
compute the maximum ungapped alignment score by scanning all start
positions and using the cumulative-sum-to-zero trick. The empirical
histogram of maxima fits the Gumbel CDF $1 - \exp(-K m n e^{-\lambda
S})$ with $K$ and $\lambda$ matching the BLAST default values
($K \approx 0.13$, $\lambda \approx 0.318$ for BLOSUM62). The right
tail matches well; the left tail is dominated by the finite-length
correction (Altschul's edge effect). A QQ-plot is the cleanest
diagnostic. This exercise verifies by simulation the analytic E-value
formula derived in the derivation section.
\end{exercise}

\begin{exercise}
Use
\repopath{volumes/06-life/ch11-bioinformatics/code/needleman_wunsch.py}
to align the pair in
\repopath{volumes/06-life/ch11-bioinformatics/data/pair_example.fasta}
under three scoring systems: $(+1, -1, -1)$, $(+2, -1, -2)$, and
$(+1, -2, -1)$. Report the optimal score and the number of gaps in
each alignment. Explain how the scoring system shifts the trade-off
between mismatches and gaps.

\paragraph{Solution sketch.}
Under $(+1, -1, -1)$, the algorithm balances mismatches and gaps
evenly; the optimal alignment has a mix of both. Under $(+2, -1, -2)$,
matches are rewarded more and gaps are penalised more; the alignment
tolerates more mismatches and fewer gaps. Under $(+1, -2, -1)$,
mismatches are penalised more than gaps; the alignment introduces gaps
to avoid mismatches. The scoring system is a prior on which
evolutionary scenario the alignment prefers; it should be chosen with
the biological context in mind, not by what makes the score look
best. This is the alignment-statistics version of choosing a prior in
Bayesian analysis: it is not optional, and it must be defended.
\end{exercise}

\subsection*{Design}\pathtag{standard}

\begin{exercise}
Design an experimental layout for an scRNA-seq study comparing tumour
and adjacent normal tissue across $40$ patients, processed in batches
of $8$ samples per sequencing run, with two sequencing centres
involved. State the per-batch composition explicitly. Justify your
choices against the failure modes in
\Cref{fig:vol06:ch11:batch-effect} and
\Cref{fig:vol06:ch11:ml-leakage}.

\paragraph{Discussion.}
Rubric. (a) Each batch contains both tumour and normal samples (no
all-tumour or all-normal batches). (b) Each batch contains samples
from multiple patients (to allow patient-level cross-validation
downstream). (c) Both sequencing centres receive a representative
mix of batches over time (no all-centre-A-early and all-centre-B-late
design that confounds time with site). (d) Patient assignment to
batches is randomised, not by enrolment order or by clinical
characteristic. (e) The block structure is documented in a metadata
sheet so that downstream analysis can use mixed-effects models or
batch-correction methods. The most defensible designs include a small
shared technical-control sample present in every batch to estimate
the residual batch shift after correction. A strong answer notes that
the design must also allow patient-level cross-validation (no patient
contributes to both training and test in any downstream ML).
\end{exercise}

\begin{exercise}
Design a benchmarking protocol for a new variant caller. Specify the
truth set, the performance metrics, the cohort selection, the
external-validation cohort, and the per-class breakdown that the
report must include.

\paragraph{Discussion.}
Rubric. (a) Truth set: Genome in a Bottle benchmarks (NA12878,
HG002) with publicly available high-confidence variant calls. (b)
Metrics: recall, precision, F1, plus per-class breakdowns (SNV vs
indel; high-confidence regions vs low-confidence regions; common vs
rare variants). (c) Cohort: at least two distinct samples processed
on at least two distinct instrument lineages. (d) External validation:
a cohort processed at an independent laboratory, ideally with a
different sequencing chemistry. (e) Per-class breakdown: separate
performance numbers for substitutions, single-base insertions,
single-base deletions, multi-base indels, and complex variants;
performance in repeats, in low-complexity regions, and at sites of
common segmental duplication, because these are the hardest classes
and the ones that distinguish callers. (f) Reporting must include
calibration: a plot of empirical accuracy versus reported QUAL
score, showing whether the caller's confidence is well calibrated.
A strong answer notes that the truth set is itself the product of a
benchmark consortium and that benchmarking against only one truth set
risks circular validation; multi-truth-set evaluation is the next
discipline up.
\end{exercise}

\begin{exercise}
Design a regulatory submission package for an omics-based diagnostic
predictor (transcriptional signature of disease subtype) intended for
clinical use. The package must satisfy the IOM 2012 recommendations.

\paragraph{Discussion.}
Rubric. (a) Locked computational predictor: every preprocessing
parameter, feature, weight, and threshold is fixed by a signed-and-dated
specification before validation begins. (b) Discovery and validation
cohorts: described in full, with inclusion and exclusion criteria,
demographics, and processing history. The validation cohort is
prospectively collected at one or more institutions independent of the
discovery cohort. (c) Predefined performance thresholds: sensitivity,
specificity, and calibration thresholds are stated in advance with the
statistical justification. (d) IRB approval that explicitly addresses
the predictor's validation status. (e) Software-as-a-medical-device
qualification per FDA pathway. (f) Post-market surveillance plan:
performance on the deployment population is monitored, and a
re-validation trigger is defined for performance drift. A strong
answer addresses the Wong et al. external-validation failure of the
Epic Sepsis Model as the cautionary parallel and explains why
post-market surveillance is not optional.
\end{exercise}

\begin{exercise}
Design a multi-omics integration plan for a study with bulk RNA-seq,
proteomics, and metabolomics measurements on $100$ patients. Specify
the per-sample sample-prep order, the storage protocol, the
identification keys, and the planned integration approach.

\paragraph{Discussion.}
Rubric. (a) Per-sample sample-prep order: simultaneous aliquoting of
fresh sample into RNA-stabilising, protein-stabilising, and
metabolite-quenching tubes at collection, because the metabolome
shifts within seconds and the proteome within minutes after sample
handling perturbations. (b) Storage protocol: $-80$~\degC{} freezer
with documented time-to-freeze for every sample. (c) Identification
keys: a single per-patient ID maps to per-tissue IDs; every aliquot
carries the patient ID, the tissue ID, the aliquot index, and the
date and time of processing. (d) Integration approach: discuss MOFA
(multi-omics factor analysis) or DIABLO for joint dimension reduction,
or a Bayesian hierarchical model that shares information across modes.
(e) Cross-modal validation: a measurement that appears in two modes
must agree to within the modes' joint uncertainty. A strong answer
addresses the central failure mode of multi-omics: a per-mode batch
effect can survive the integration if it is consistent across modes,
because the integration estimator may interpret it as a real
multi-omics signal.
\end{exercise}

\subsection*{Diagnosis}\pathtag{standard}

\begin{exercise}
A colleague reports that their RNA-seq differential-expression
analysis identified $10{,}000$ genes at FDR $< 0.05$, comparing
$8$ cancer samples against $8$ normal samples. They are excited and
write a draft paper. What questions do you ask before allowing the
analysis into the paper?

\paragraph{Solution sketch.}
(1) How are the $16$ samples blocked across batches? (Look for
batch-condition collinearity.) (2) What is the per-sample mapping
rate and the per-sample ribosomal-read fraction? (Look for sample
outliers that drive the result.) (3) What is the dispersion estimate
for the typical gene? (A $10{,}000$-gene DE list is biologically
implausible unless the conditions are very different; a much more
common cause is under-shrinkage of the dispersion estimate due to
batch effect or sample contamination.) (4) Is the patient and the
sample collection time recorded for each sample? (Same-patient
matched-pair design changes the right test.) (5) Plot PCA before
and after batch correction. (6) Plot the volcano plot and examine
whether the top hits are biologically sensible or appear to be
arbitrary high-variance features. (7) Repeat the analysis with one
sample removed at a time and check the stability of the gene list.
The most common answer to "$10{,}000$ DE genes at FDR $< 0.05$" in
a $16$-sample comparison is "batch effect dominates the analysis."
\end{exercise}

\begin{exercise}
A protein-structure prediction returns a 600-residue model with
average pLDDT $52$ and one well-folded domain of pLDDT $85$ at
residues 200--350. Diagnose what this prediction does and does not
support, and recommend follow-up.

\paragraph{Solution sketch.}
Mean pLDDT of $52$ indicates that the bulk of the model is at low
confidence and probably does not correspond to a single experimentally
observable conformation. The high-pLDDT domain at 200--350 is at
near-experimental accuracy for the local structure and is a credible
prediction of the domain fold. The model supports: a domain-level
fold call for residues 200--350; the inter-domain topology only
weakly. The model does not support: atomic-coordinate-level claims
about residues outside 200--350; a single biological conformation if
the low-pLDDT regions are intrinsically disordered. Recommended
follow-up: check the IUPred or other disorder-prediction tools for
the low-pLDDT regions; if they predict disorder, the AlphaFold result
is consistent and the protein has one folded domain plus disordered
arms. If they do not predict disorder, consider that the low-pLDDT
regions may fold but in a state under-represented in training data
(membrane embedded, multimerised, ligand bound); recommend
ESMFold or AlphaFold-Multimer with the appropriate partners.
\end{exercise}

\begin{exercise}
A clinical variant report identifies a variant of uncertain
significance (VUS) at a known disease gene. The variant has allele
frequency $0.001$ in gnomAD, is predicted deleterious by SIFT and
PolyPhen, and is in a conserved residue. The variant is absent from
ClinVar. Diagnose what the report can support clinically.

\paragraph{Solution sketch.}
Allele frequency $0.001$ is rare but not vanishingly rare; the
variant has been seen $\sim 100$ times in gnomAD's $\sim 10^{5}$
exomes, which is too common for a fully penetrant Mendelian
disease and too rare to dismiss. SIFT and PolyPhen are computational
predictors with imperfect specificity; a "deleterious" call from
both is suggestive but not diagnostic. Conservation supports
functional importance but does not establish a disease link. Absent
from ClinVar means that no clinical-significance call has been
issued, not that the variant is benign. The ACMG/AMP variant-
interpretation framework would weight these as moderate computational
evidence (PP3), moderate population evidence (PM2 absent or rare),
and would not assign pathogenic without additional segregation,
functional, or phenotypic evidence. The report supports: continued
clinical follow-up; family-segregation testing if the patient has an
affected relative; functional assay if the gene's biology permits
one. The report does not support: a confident pathogenic call. The
correct clinical action is to communicate the variant as a VUS and to
defer treatment decisions to the gene-and-disease-specific
guidelines.
\end{exercise}

\subsection*{Judgement}\pathtag{mastery}

\begin{exercise}
Compare the ESMFold and AlphaFold approaches to protein structure
prediction. Under what conditions would you trust ESMFold and not
AlphaFold? Under what conditions would you trust AlphaFold and not
ESMFold? Construct a one-paragraph decision rule for a working
structural biologist.

\paragraph{Discussion.}
ESMFold uses a protein language model trained on hundreds of millions
of unaligned sequences and predicts structure from a single sequence.
AlphaFold uses a multiple-sequence alignment of orthologues and a
template search of the PDB. AlphaFold's accuracy depends on the
quality of the MSA; for proteins with deep MSAs (most well-studied
families), it is the more accurate predictor by a sizeable margin.
For proteins with no or very few orthologues (orphans, designed
proteins, fast-evolving viral proteins, synthetic constructs),
AlphaFold has nothing to work with and ESMFold's single-sequence
prediction is the better tool. ESMFold is also two orders of
magnitude faster, making it preferred for proteome-scale screens
where the goal is to identify high-confidence targets for follow-up
with AlphaFold. The working decision rule: if the protein has $\geq
50$ diverse orthologues in UniRef, run AlphaFold and read the pLDDT;
if it has $< 10$ orthologues or is a designed sequence, run ESMFold
first and confirm with AlphaFold if the ESMFold prediction is
biologically credible. A strong answer notes that both architectures
share a common limitation: they predict the dominant conformation in
the training distribution, not the relevant one for a particular
biological question (membrane, partner-bound, allosteric state).
\end{exercise}

\begin{exercise}
A genomics start-up claims that its AI-based predictor identifies
patients who will respond to a novel chemotherapy with sensitivity
$95\%$ and specificity $90\%$, on the basis of a discovery cohort of
$240$ patients with $5$-fold cross-validation. The predictor's
mechanism is proprietary. The company offers the predictor to your
clinical-trials office for use in patient selection. Evaluate the
offer.

\paragraph{Discussion.}
The offer fails the IOM 2012 recommendations on at least three counts.
(1) No external-validation cohort is reported; cross-validation on
the discovery cohort is a development metric, not a release metric.
(2) The predictor's mechanism is proprietary; the IOM framework
requires locked computational predictors with explicit documentation
of the algorithm, the features, and the training data. (3) The cohort
size is small ($240$ patients) and the predictor's apparent
performance is implausibly high for chemotherapy response, which is
notoriously multi-factorial. The correct action is to decline the
offer pending: an independent external-validation cohort processed at
the trials-office institution; a release of the predictor's algorithm
and training set (under NDA acceptable) for independent reanalysis;
and IRB approval that explicitly addresses the predictor's validation
status. A strong answer cites Duke Potti as the precedent for
declining this kind of offer and the consequences of accepting it.
\end{exercise}

\begin{exercise}
A drug-target predictor based on a graph neural network achieves
$0.92$ AUC on a held-out test set drawn from the same dataset as the
training set. When applied prospectively to identify candidate
targets for follow-up screening, $2$ of $50$ predicted targets
validate experimentally. Is the model good, bad, or unevaluable from
this evidence?

\paragraph{Discussion.}
$0.92$ AUC on within-distribution test and $2/50 = 4\%$ prospective
hit rate is consistent with three explanations. (1) Same-distribution
overfit: the test set shares structure with the training set
(scaffold leakage, protein-family leakage) and the AUC is optimistic.
(2) Distribution shift: the prospective targets come from a different
chemical space than the training data; the model's predictions are
unreliable out of distribution. (3) The prospective screen is the
correct evaluation: $4\%$ hit rate is in fact much better than a
random pick from the prior list ($\sim 1\%$ baseline for most drug-
target screens), and the model is operating as designed. Without more
information, the evidence is consistent with all three. The
diagnostic action is to check the scaffold-clustering and protein-
family-clustering structure between training set and prospective
candidates; if there is no overlap, the model is doing useful work at
the level its evidence supports. A strong answer notes that the
relevant comparison is not against the within-distribution AUC, which
is a development metric, but against a random-baseline prospective
screen, which is the deployment metric. This is the same
internal-versus-external evaluation distinction that anchors the
failure section.
\end{exercise}

\subsection*{Failure analysis}\pathtag{mastery}

\begin{exercise}
Reconstruct the data-leakage mechanism in the Duke Potti case from the
short-form summary in
\repopath{docs/research/accidents/duke-potti-2006-2015.md}. For each
of the four canonical leakage pathways in
\Cref{fig:vol06:ch11:ml-leakage}, identify whether the pathway
contributed and explain. Propose three structural changes to the
NIH-funded translational-omics pipeline that would have caught the
errors before clinical-trial enrolment.

\paragraph{Solution sketch.}
Pathway 1 (sample-level versus patient-level split): contributed
partly. The predictor was trained on NCI-60 cell lines, not on
patients, so the conventional split was at the cell-line level. The
specific error Baggerly and Coombes identified was that some cell
lines appeared in both training and test sets due to off-by-one
shifts. Pathway 2 (preprocessing on full cohort): contributed.
Feature-selection rankings reported in the original papers used the
test set's labels in their selection, leaking outcome information
into the training pipeline. Pathway 3 (hyperparameter tuning on test
set): contributed. The reported performance numbers were optimistic
because thresholds were tuned on the same data on which they were
evaluated. Pathway 4 (same-source contamination): not the dominant
mechanism; the NCI-60 cell lines have known same-source structure but
that was understood. Three structural changes to the NIH pipeline:
(a) require a locked computational predictor (per IOM 2012
Recommendation 1) before validation begins, with the full code and
data deposited and reproducibly re-executable; (b) require independent
reanalysis by a designated group (the role Baggerly-Coombes played
unofficially) before any clinical-trial enrolment can begin; (c)
require that the IRB application for any clinical trial using an
omics predictor explicitly state the validation status of the
predictor and the criteria by which it was judged ready for clinical
use. A strong answer notes that the family-of-failures pattern (Duke
Potti, COVID-19 imaging, Epic Sepsis) is the canonical evidence for
the post-2012 IOM framework being the right minimum standard.
\end{exercise}

\begin{exercise}
A research paper reports a deep-learning model that predicts
pancreatic-cancer survival from H\&E-stained histology slides with
$0.83$ AUC on a $400$-patient cohort with $5$-fold cross-validation.
The cohort comes from a single institution. The training and test
folds were stratified by patient at the slide-image level. No
external validation was performed. The paper makes no claim about
clinical use; it is published as a methods demonstration. Evaluate
the paper's claims and propose what it should be required to add
before clinical deployment.

\paragraph{Solution sketch.}
The paper's claim (a methods demonstration on a $400$-patient cohort
with patient-level cross-validation) is supported by the evidence
presented. The reported AUC is on a within-institution distribution
and is a development metric. The paper does not claim clinical
utility; the methods-demonstration framing is appropriate. For
clinical deployment, the paper must add: an external-validation
cohort from at least one independent institution (preferably more,
to span scanner and staining variation); a calibration analysis (the
posterior probability of survival at the predicted threshold should
match the empirical survival); a per-stratum performance breakdown
(by stage, by treatment, by demographic) to identify subgroups where
the model fails; a sensitivity analysis to known confounders
(scanner instrument, staining lot, slide preparation time, patient
demographics); and a pre-registered analysis plan for the external
validation. The Roberts et al. COVID-19 review is the canonical
reference for the failure mode the external-validation analysis
guards against: $415$ comparable papers exist and none are clinically
usable. The paper's authors have done their part as a methods
demonstration; the path from methods to deployment requires another
study, not a re-interpretation of this one. A strong answer
distinguishes between two release tiers: the methods-demonstration
tier (this paper) and the clinical-validation tier (a separate study
yet to be done), and refuses to conflate them.
\end{exercise}
